\documentclass[11pt,reqno,a4paper]{amsart}
\usepackage[T1]{fontenc}
\usepackage{lmodern,mathtools,amssymb,amsthm,booktabs,array,enumitem}
\usepackage[expansion=false]{microtype}
\usepackage[a4paper,textwidth=153mm,textheight=232mm,centering]{geometry}
\usepackage[hidelinks]{hyperref}
\hypersetup{pdftitle={Full index prefixes in combinatorial sequences and simultaneous digit windows},pdfauthor={Nikita Maximilian Grimm},pdfkeywords={radix expansion, Subspace Theorem, Hausdorff dimension, open dynamical system}}
\pdfinfoomitdate=1
\pdftrailerid{}
\pdfsuppressptexinfo=15
\newtheorem{theorem}{Theorem}[section]
\newtheorem{lemma}[theorem]{Lemma}
\newtheorem{proposition}[theorem]{Proposition}
\newtheorem{corollary}[theorem]{Corollary}
\theoremstyle{remark}
\newtheorem{remark}[theorem]{Remark}
\newtheorem{example}[theorem]{Example}
\newcommand{\companion}[1]{\href{../#1.md}{\nolinkurl{#1.md}}}
\setlength{\emergencystretch}{1.2em}
\raggedbottom
\setlist[enumerate,1]{leftmargin=*,label=\textup{(\roman*)}}
\title[Full index prefixes and digit windows]{Full index prefixes in combinatorial sequences\\ and simultaneous digit windows}
\thanks{Curated by Nikita Maximilian Grimm with AI assistance. Repository: \url{https://github.com/NikitaMGrimm/math}.}
\date{}
\begin{document}
\begin{abstract}
For every fixed pair of integers \(s,t\), the binary expansion of
\(\binom{2n+s}{n+t}\) begins with the entire expansion of \(n\) at a number
of indices in \([X,X^{3/2+\varepsilon}]\) tending to infinity. An
analogous result holds for centered, strongly aperiodic planar lattice
walks with a power of two labeled steps. Up to \(X\) both families have
at least a constant times \((\log\log X)(\log\log\log X)\) matches.
For counts asymptotic to
\(Cb^{hn}n^{1-p}(1-\gamma/n+O(n^{-2}))\), we classify the integer
radices admitting a zero-entropy survivor for the associated circle
map. There are exactly seven, each realized by an unshifted count.
The arithmetic argument also yields exact Hausdorff dimensions for long
simultaneous zero blocks on an algebraic line. Synchronizing two
multiplicatively independent radices strictly lowers the critical block
exponent. For integer-valued polynomials we decide when the index occurs
as both prefix and suffix infinitely often. In Zeckendorf notation the
large indices prefixing \(n^2+c\) are exactly the Lucas numbers \(L_k\)
with \(c+2(-1)^k\ge0\).
\end{abstract}
\maketitle

\section{Introduction}

The decimal expansion of \(2^{10}=1024\) begins with the index \(10\).
Erd\H{o}s asked whether this happens at infinitely many indices; see
van de Lune~\cite{vdl} and Fang~\cite{fang}. The target changes with
the index. Prescribing the entire expansion of an integer of size \(n\)
amounts to prescribing an interval of relative width \(1/n\).
In this paper a polynomial factor in an exponential count makes it
possible to convert this shrinking interval into a fixed arc on the
circle. We first give the resulting prefix theorems, then classify the
arcs used by the proof and determine the size of the exceptional sets
in the underlying problem on digit expansions.

For a positive integer sequence \(U=(U_n)\), an integer radix \(b\ge2\),
and \(X\le Y\), put
\[
 H_b(U;X,Y)=\#\{n\in\mathbb N\cap[X,Y]:
             nb^k\le U_n<(n+1)b^k\ \text{for some integer }k\ge0\}.
\]
The block of \(k\) digits following the prefix may begin with zeros. Write \(W_b(U;\sigma)\) if
\(H_b(U;X,X^{\sigma+\varepsilon})\to\infty\) as \(X\to\infty\)
for every \(\varepsilon>0\). Set
\[
 B_n=\binom{2n}{n},\qquad
 T_n=[z^n](1+z+z^2)^n,\qquad F_r(n)=\sum_{j=0}^n\binom nj^r.
\]

\begin{theorem}[Combinatorial prefixes]\label{thm:headline}
\leavevmode\begin{enumerate}
\item For every \(s,t\in\mathbb Z\),
 \(W_2(\binom{2n+s}{n+t};3/2)\) holds on the eventually defined sequence.
 For \(r=3,5,7\) the assertions \(W_2(F_r;2)\),
 \(W_2(F_r;3/2)\), \(W_2(F_r;4/3)\) hold, respectively.
 For integer \(r\ge2\) the corresponding binary survivor has positive
 entropy exactly at even orders \(r\ge4\) and odd orders \(r\ge9\).
\item Let a centered list of \(2^h\) lattice steps in the plane have
 pairwise differences generating its rank-two lattice of positions \(\Lambda\), with \(h\ge1\).
 For every fixed \(z\in\Lambda\), the count \(U_n(z)\) of
 \(n\)-step walks from \(0\) to \(z\) satisfies \(W_2(U(z);2)\).
 Repeated steps count as distinct choices. In ternary, \(W_3(T_n^2;2)\) holds.
\item There are infinitely many full prefixes for \(B_n^2\) in base four,
 \(F_3(n)^2\) in base eight, \(T_n^2\) in base nine, and \(F_4(n)^2\)
 in each of bases two, four and sixteen. In base \(32\) there are
 infinitely many for the following return count on \(\mathbb Z^8\).
 Arrange the coordinate labels on a disjoint triangle and pentagon;
 choose an edge and change both its coordinates by independent signs.
\end{enumerate}
\end{theorem}

For example, \(B_4=70=(1000110)_2\) begins with \(4=(100)_2\), and \(2^{-27}B_{17}\in[17,18)\). To see what selects these indices, suppose that \(B_n/2^k=n+f\), \(0\le f<1\). Stirling's formula gives
\[
 \pi^{-1/3}4^{(2n-k)/3}=n+\tfrac1{12}+\tfrac23 f+O(1/n).
\]
Thus one tests a fixed interval with three interlaced geometric
sequences. Their common factor \(\pi^{-1/3}\) need not satisfy any
Diophantine hypothesis for Theorem~\ref{thm:headline}(i).
Theorem~\ref{thm:app-counting} gives the cumulative lower bound for
these three branches and the other multibranch examples.

For the general construction suppose that
\begin{equation}\label{eq:aligned}
 U_n=Cb^{hn}n^{1-p}\bigl(1-\gamma/n+O(n^{-2})\bigr),
 \qquad C>0,\ h\in\mathbb N,\quad p=P/Q>1
\end{equation}
with \(P,Q\) coprime positive integers. Factor \(b=g^u\) with \(g\)
an integer which is not a perfect power. Define
\begin{equation}\label{eq:branch}
 l=\gcd(P,u),\qquad e=P/l,\qquad \tau=b^{1/p},\qquad B=\tau^e=g^{uQ/l}.
\end{equation}
On \(\mathbb R/\mathbb Z\) take \(T_B(x)=Bx\bmod1\).
For an open arc \(I\) write \(K_B(I)=\{x:T_B^m x\notin I\ (m\ge0)\}\).
We call it the survivor and use the topological entropy of \(T_B\)
on this compact forward-invariant set (zero for the empty set).
Our arc is
\begin{equation}\label{eq:hole}
 I=\bigl(\gamma/p,(\gamma+1)/p\bigr)\pmod1
\end{equation}
with positive orientation and length \(1/p\).

\begin{theorem}[Radix classification and transfer]\label{thm:transfer}
For every real \(\gamma\) the survivor \(K_B(I)\) has zero entropy
exactly in the cases of Table~\ref{tab:zones}. As \(p\) and \(\gamma\)
vary, the possible numeral radices are precisely
\[
                  2,\ 3,\ 4,\ 8,\ 9,\ 16,\ 32 .
\]
If \(e>1\), \(\gamma\) is rational, and \(K_B(I)\) has zero entropy,
then \(W_b(U;e/(e-1))\) holds for every sequence satisfying
\eqref{eq:aligned}. The same implication holds for real \(\gamma\)
if an open arc \(J\) with rational endpoints, \(\overline J\subset I\),
has zero-entropy survivor.

If \(e=1\), assume such an inner arc \(J\) exists and
\(\alpha=C^{1/p}\) is irrational. There are infinitely many full
matches. If \(\alpha\) has finite irrationality exponent \(\mu\), let
\(s\) be the maximum number of cyclic components visited by a path in
the finite graph for \(J\) defined in Section~\ref{sec:transfer}.
Every sufficiently large \([X,X^\rho]\) contains a match for
\(\rho>\max(1,\mu^s)\).
\end{theorem}

\begin{table}[ht]
\centering
\caption{Complete zero-entropy parameters for the branch multiplier \(B\).
All interval endpoints are included. Each row uses its displayed period
in \(\gamma\).}\label{tab:zones}
\begin{tabular}{ccl}
\toprule
\(B\)& \(p\)& \(\gamma\)\\
\midrule
\(2\)&\(2\)& all real values\\
\(2\)&\(3\)& \([1/2,3/2]\pmod 3\)\\
\(2\)&\(4\)& \([1,2]\pmod 4\)\\
\(2\)&\(5\)& \([23/12,25/12]\pmod 5\)\\
\(3\)&\(2\)& \([1/3,2/3]\pmod 1\)\\
\(4\)&\(3/2\)& \([1/8,3/8]\pmod{1/2}\)\\
\bottomrule
\end{tabular}
\end{table}

The classification concerns the one-interval criterion in
\eqref{eq:branch}--\eqref{eq:hole}. Positive survivor entropy leaves
the behavior of a specific sequence open. The doubling-map part uses the
asymmetric-hole theorem of Glendinning and Sidorov; its specialization
to the widths \(1/p\) is detailed in Section~\ref{sec:classification}.

The arithmetic problem behind the transfer also has a metric form.
Let \(v=(v_1,\ldots,v_d)\) consist of positive real algebraic numbers
linearly independent over \(\mathbb Q\), with \(d\ge2\). Use digits after
the radix point, choosing the terminating expansion when available.
For a base \(B_1\), a start \(u\in\mathbb N\), and \(\rho>1\), a zero
window for a scale \(c>0\) means that digits \(u+1,\ldots,\lfloor\rho u\rfloor\)
of every \(cv_i\) vanish. In a second base \(B_2\) synchronize the
start \(u_2=\lfloor u\log_{B_2}B_1\rfloor\) and end at
\(\lfloor\rho u_2\rfloor\). Denote the sets of multipliers satisfying
these conditions for infinitely many \(u\) by \(E_1(\rho)\) in one
base and \(E_{12}(\rho)\) in both. Write \(x_+=\max(x,0)\).

\begin{theorem}[Simultaneous digit windows]\label{thm:dimension}
The Hausdorff dimensions in every nonempty open subinterval of
\((0,\infty)\) are
\begin{align*}
 \dim_H E_1(\rho)&=\bigl(d/\rho-(d-1)\bigr)_+,\\
 \dim_H E_{12}(\rho)&=\rho^{-1}\max\{0,\,2d-(2d-1)\rho,\,
                                      d\lambda-(d-1)\rho\},\\
 \lambda&=\sum_{\ell\ {\rm prime}}
     \min\left\{\frac{\nu_\ell(B_1)\log\ell}{\log B_1},
                 \frac{\nu_\ell(B_2)\log\ell}{\log B_2}\right\}.
\end{align*}
The sets are empty, respectively, above \(d/(d-1)\) and above
\(\max\{2d/(2d-1),d\lambda/(d-1)\}\). Every multiplicatively
independent pair strictly lowers the critical exponent.
At the one-base endpoint \(d/(d-1)\), if the \(v_i\) form a basis
of a degree-\(d\) number field under a real embedding, the exceptional
set contains a dense \(G_\delta\) on some half-line.

For \(1<\rho<3/2\), the positive \(g\) for which
\(H_2(\lfloor gB_n\rfloor;X,X^\rho)=0\) at arbitrarily large \(X\)
have local dimension \(3/\rho-2\). The same dimension holds if one
requires only that these counts fail to tend to infinity. At \(\rho=3/2\)
both dimensions are zero, and the empty-window parameters contain a
dense \(G_\delta\) in \((0,\infty)\). Above \(3/2\) both sets are empty.
\end{theorem}

For instance, for the line \(c(1,\sqrt2)\) the cutoff is \(2\) in
binary alone and \(4/3\) in binary and ternary together. The parameter
\(\lambda\) measures exactly the size of the common denominator grid.
When \(N=B_1^u\) and \(M=B_2^{u_2}\), one has \(\gcd(N,M)\asymp N^\lambda\).
For comparison, still in dimension two, take \(B_1=2^{10}\) and
\(B_2=3\cdot 2^{10}\). Here \(\lambda=\log(2^{10})/\log(3\cdot 2^{10})>2/3\).
Their joint cutoff is \(2\lambda\), strictly between \(4/3\) and \(2\).

Section~\ref{sec:polynomials} also tests the last digits of a polynomial.
For example, \(n^2-n+1\) begins and ends with \(n\) at \(n=b^m+1\)
for every sufficiently large \(m\), whereas \(n^d+c\) does this only
finitely often for each fixed \(b,d\ge2\) and integer \(c\).
Section~\ref{sec:representations} treats changes to the digit word
itself, including reversal and the Zeckendorf representation.

\section{From leading digits to orbits on the circle}\label{sec:transfer}

Write \(\alpha=C^{1/p}\). For a trial suffix \(k=hn-j\), the natural
continuous index is \(x=\alpha\tau^j\). If \(x-n\) is bounded, then
\eqref{eq:aligned} gives, uniformly in this range as \(x\to\infty\),
\begin{equation}\label{eq:residual}
 b^{-k}U_n-n=p(x-n)-\gamma+O(x^{-1}).
\end{equation}
Indeed \(Cb^j=x^p\), and expansion of
\(n(x/n)^p(1-\gamma/n+O(n^{-2}))\) gives the formula.
For \(n=\lfloor x-\gamma/p\rfloor\), the right side is
\(p\{x-\gamma/p\}+O(x^{-1})\); also \(k\) is eventually positive.
A phase in a compact subarc of \(I\) therefore gives a match. For
converse arguments, a genuine match with \(b^{-k}U_n=n+f,\ 0\le f<1\),
first gives \(x/n=1+O(1/n)\) by taking a \(p\)-th root, and then
\begin{equation}\label{eq:inverse}
 x=n+(\gamma+f)/p+O(n^{-1}).
\end{equation}
The latter inclusion need not select the same floor at an endpoint.

The positive number \(\tau\) in \eqref{eq:branch} has degree exactly
\(e\). If its degree is \(d'\), norms in \(\tau^e=B\) give
\(e\mid d'\nu_\ell(B)\) for each prime \(\ell\mid g\).
The gcd of the valuations of \(g\) is one and
\(\gcd(e,uQ/l)=1\), so \(e\mid d'\). The equation \(T^e-B=0\)
gives \(d'\le e\). Thus in the branches \(j=em+r\), \(0\le r<e\),
the \(x\)'s at time \(m\) are the coordinates of
\begin{equation}\label{eq:line}
             \alpha B^m(1,\tau,\ldots,\tau^{e-1}).
\end{equation}

\subsection{An arithmetic bound for sparse transitions}
Digit-change estimates obtained by the Subspace Theorem ordinarily
concern an individual algebraic irrational; see, for example,
Bugeaud and Evertse~\cite{be-digits}. In the lemma below the common
multiplier may be any nonzero real number, and the approximants need
only stay at bounded distance from the line.

\begin{lemma}\label{lem:transitions}
Let \(B\ge2,\ d\ge2\) be integers, let \(v_1,\ldots,v_d\) be
positive real algebraic numbers linearly independent over \(\mathbb Q\),
and fix \(c\ne0\). If \(W_m\in\mathbb Z^d\) satisfies
\(W_m=cB^m v+O(1)\) coordinatewise, then for every \(\eta>d/(d-1)\),
\[
 \#\{u<m\le\lfloor\eta u\rfloor:W_m\ne BW_{m-1}\}\longrightarrow\infty .
\]
\end{lemma}
\begin{proof}
Suppose the counts stay bounded along a subsequence and let
\(w=\lfloor\eta u\rfloor\). Each transition vector \(W_m-BW_{m-1}\)
belongs to a fixed finite set. Pass to a subsequence fixing the
number and ordered values of the nonzero transitions. Each gap
between successive transition times can in turn be held constant
or made to tend to infinity. Group times separated by constant
gaps and write \(a_i\) for the last time of each resulting group.
Telescoping yields
\[
 W_w=Y+\sum_{i=1}^s h_i P_i,\qquad
 Y=B^{w-u}W_u,\qquad P_i=B^{w-a_i}.
\]
The vectors \(h_i\) are fixed integers: contributions from earlier
times of a group are multiplied by fixed nonnegative powers of \(B\).
Discard zero group contributions. The remaining \(a_i\)'s are ordered
with pairwise gaps tending to infinity.

We use the \(S\)-adic Subspace Theorem in its integer-product form
\cite{schlickewei}. For \(m\ge2\) independent algebraic linear forms
in \(m\) variables at each place of a finite set \(S\) containing
the real place, integer vectors with
\(\prod_{\nu\in S}\prod_i |L_{i,\nu}(Z)|_\nu
 \le \|Z\|_\infty^{-\delta}\), \(\delta>0\), lie in finitely many
proper rational subspaces. The finite absolute values are normalized
by \(|\ell|_\ell=\ell^{-1}\). In this proof the finite-place forms
are the coordinate forms. The version for primitive vectors implies
this one by dividing out the common gcd (its product over \(S\) is
at least one).

Take \(Z=(Y_1,\ldots,Y_d,P_1,\ldots,P_s)\), with finite places the
prime factors of \(B\). At infinity put \(F=Y+\sum h_i P_i\) as
a vector of forms, and choose
\[
 F_1,\quad F_r-(v_r/v_1)F_1\quad(2\le r\le d),
     \qquad P_1,\ldots,P_s .
\]
These forms are independent. The differences evaluated at \(Z\)
are \(O(1)\); \(F_1=O(B^w)\). Each \(Y_r\) contributes at most
\(B^{-(w-u)}\) at the finite places, and each \(P_i\) has product
one over all selected places. Thus the full product is
\(O(B^{du-(d-1)w})\). Since \(\|Z\|_\infty\asymp B^w\), this is a
negative power of the height. An infinite subsequence satisfies
one fixed nontrivial rational linear relation. Divide it by \(B^w\).
The \(Y\) terms tend to \(cv\), and all \(P_i/B^w\) tend to zero.
Independence of the \(v_r\)'s kills every coefficient of \(Y\).
Among the remaining powers, divide by the largest with nonzero
coefficient. All smaller powers disappear in the limit, a contradiction.
\end{proof}

\subsection{The endpoint error and the finite graph}

For an open rational arc \(J\), take a common denominator \(D\) of
its endpoints and partition the circle into the half-open cells
\([i/D,(i+1)/D)\), \(0\le i<D\). Keep those whose interiors are
outside \(J\). Draw the edges
\[
             i\longrightarrow Bi+t\bmod D,\qquad 0\le t<B,
\]
between retained states, retaining multiplicities and labels \(t\).
Indeed for \(V_m=\lfloor D x_m\rfloor\), \(x_{m+1}=Bx_m\),
one has \(V_{m+1}=BV_m+t_m\), where \(t_m\) is the next radix digit
of \(\{D x_m\}\). A cyclic strongly connected component is a component
containing a directed cycle.

This graph has zero entropy exactly when each cyclic component is a
simple directed cycle. A branching component admits two distinct return
words from a vertex with different initial edges, which can be repeated
to give two distinct equal-length return words. Simple cyclic components
give polynomially many paths of each length, since a path can change
components only boundedly often. This criterion gives the entropy-zero
test for \(K_B(J)\). Indeed an infinite graph path gives a point of the
closed survivor by intersecting its nested closed itinerary intervals.
For positive-entropy paths the coding by the digits of \(Dx\) identifies
at most the two expansions at a radix endpoint for a fixed initial cell.
Conversely a survivor missing the grid visits only retained cells. Points
which ever visit the grid form a countable set of eventually periodic
points. More formally one may code the closed survivor in the full
one-sided radix shift, using both codes of a terminating expansion.
The coding is a finite-to-one factor and preserves entropy (the fibers
carry zero relative entropy). An ergodic positive-entropy measure gives
no mass to the countable grid preimages. The variational principle
therefore identifies positivity on either side, including holes which
wrap around zero.

Suppose first \(e>1\) and \(I\) itself is rational with a zero-entropy
graph. Work on one of the \(e\) branches \(x_m=\alpha B^m\tau^r\),
on \([u,w]\) with \(w=\lfloor\eta u\rfloor\) and fixed
\(\eta>e/(e-1)\). At a failed test in a nonretained cell,
\eqref{eq:residual} gives distance \(O(B^{-m})\) to a grid endpoint,
since phases in the interior of the success arc farther than this
quantity from both ends give a true hit. This includes a grid boundary
assigned to a nonretained cell. If such a failure occurs at \(a\),
\(\|D x_a\|_{\mathbb R/\mathbb Z}\ll B^{-a}\). For
\(a\le m\le\min(2a-K,w)\), the nearest integer to \(D x_m\)
is \(B^{m-a}\) times that at \(a\); choose \(K\) so that the propagated
error stays strictly below \(1/2\). Cover these failures in increasing
order, beginning a new interval at the first uncovered one. New starts
are at least \(2a-K+1\). Thus \(O_\eta(1)\) intervals suffice.

Delete those intervals and all successful tests for the moment. If at
most \(H\) tests succeeded, the remaining \(O_\eta(H+1)\) runs follow
retained paths. On each stretch within a cyclic component let
\(\theta_i\in[0,1]\) be the radix fraction read by continuing around
its unique cycle from state \(i\). An edge \(i\to i'\) with digit
\(t\) satisfies \(\theta_{i'}=B\theta_i-t\), also for the constant
word of digits \(B-1\), which gives \(\theta_i=1\).
Let \(A\) be a common multiple of the \(B^q-1\), \(1\le q\le D\).
Use the integer \(A(V_m+\theta_i)\) at a cyclic state on a retained
run. It multiplies exactly by \(B\) on each internal cycle edge.
Use \(A\) times the nearest integer to \(D x_m\) everywhere else.
Within the deleted endpoint intervals those integers also multiply by
\(B\). There are at most \(D-1\) component changes per retained run.
All these approximants have error \(O(1)\) from \(ADx_m\);
the number of transitions on which multiplication can fail is
\(O_{\eta,D}(H+1)\).

If the total successes stayed bounded along infinitely many such
windows, take a disjoint subsequence of windows and use these integer
approximants on all branches there, filling the gaps by rounding.
This contradicts Lemma~\ref{lem:transitions} for
\(AD\alpha(1,\tau,\ldots,\tau^{e-1})\), with the scalar \(AD\alpha\)
outside the algebraic direction. For a target exponent
\(e/(e-1)+\varepsilon\) choose \(\eta\) strictly between the two
exponents, and take \(u=\lceil\log_B(X/\alpha)\rceil+O(1)\) with
a sufficiently large upward offset. All candidate indices lie in the
required interval and are eventually distinct because consecutive
unrounded \(x\)'s grow by the fixed ratio \(\tau>1\).
If a rational inner arc \(J\) is used in place of \(I\), every point
in its closure succeeds eventually, and the same argument applies
without exceptional intervals. This proves the multibranch transfer.

Finally suppose \(e=1\) and use the rational inner arc. If \(\alpha\)
is irrational, its orbit avoids rational grid points. An infinite
run of failures must eventually traverse a single periodic cycle.
The base-\(B\) digits of \(D\alpha\) would then be eventually periodic,
which is impossible. For a quantitative estimate, a stretch of \(L\)
cycle edges starting at time \(m\) gives an approximation to
\(\alpha\) of denominator \(O(B^m)\) and error \(O(B^{-m-L})\)
by continuing the edge digits periodically. For every \(\delta>0\)
the definition of \(\mu=\mu(\alpha)\) implies
\(m+L\le(\mu+\delta)m+O(1)\); bounded reduced denominators also
satisfy the needed error lower bound after adjustment of its constant.
Off the cycles at most \(O(D)\) edges occur in a retained run.
Iterating over at most \(s\) cycles and leaving a strict exponent
margin proves the last assertion of Theorem~\ref{thm:transfer}.

\section{The survivor classification}\label{sec:classification}

Write \(T_B x=Bx\bmod 1\) on \(\mathbb R/\mathbb Z\). Let \(I\) be the
open arc from \(\gamma/p\) to \((\gamma+1)/p\) in the positive
direction, \(J\) its closed complement, and
\(K=\bigcap_{j\ge0}T_B^{-j}J\). Put \(w=1/p\) and \(L=1-w\).
For a closed invariant set of \(T_B\), its Hausdorff dimension is
\(h_{\rm top}/\log B\). This is the constant-derivative case of the
pressure formula in \cite[Theorem~3]{app-urbanski}. We will use dimension
and entropy interchangeably when testing positivity.

Translation of \(\gamma\) by \(p\) gives the same hole. Reflection of the
circle sends \(\gamma\) to \(-\gamma-1\). Translation of the circle by a
multiple of \(1/(B-1)\) also commutes with \(T_B\), so it gives a period
\(p/(B-1)\) for the entropy as a function of \(\gamma\).

\subsection{An inverse-image bound}

If \(BL\ge2\), every point has at least \(\lfloor BL\rfloor\) inverse
images in \(J\), since its \(B\) inverse images are spaced by \(1/B\).
The arc \(J\) contains a fixed point because
\(L\ge 2/B\ge1/(B-1)\). Taking the inverse tree of such a point inside
\(J\) gives at least \(\lfloor BL\rfloor^n\) distinct points of \(K\)
at depth \(n\). For any two of these points, at the last time before
their orbits meet they are distinct inverse images of the same point
and are separated by at least \(1/B\). Thus
\begin{equation}\label{eq:app-inverse-tree}
 h_{\rm top}(T_B|_K)\ge \log\lfloor BL\rfloor \qquad(BL\ge2).
\end{equation}

For an admissible branch multiplier, write \(p=P/Q\) in lowest terms,
\(b=g^u\) with \(g\) not a perfect power, and \(\ell=\gcd(P,u)\), so that
\(B=g^{uQ/\ell}\). In particular \(B\) is a \(Q\)-th power. For \(Q\ge3\),
\[
 BL\ \ge\ \frac{2^Q}{Q+1}\ \ge\ 2.
\]
For \(Q=2\), we have \(P\ge3\) odd. If \(P\ge5\), then
\(BL\ge4(1-2/5)>2\). If \(P=3\) and \(B\ge9\), then \(BL\ge3\).
This leaves \((B,p)=(4,3/2)\). For \(Q=1\), \eqref{eq:app-inverse-tree}
applies to \(B\ge4\) and also to \(B=3,\ p\ge3\). The remaining
integer-\(p\) cases are tripling at \(p=2\) and doubling.

\subsection{Tripling and the half-integral exception}

For \((B,p)=(3,2)\) and \((4,3/2)\) we have \(L=1/(B-1)\).
The period in \(\gamma\) given by circle translation is \(pL\),
which is \(1\) and \(1/2\), respectively. Reflection is
\(\gamma\mapsto-\gamma\) modulo \(pL\) in both cases. We may therefore
take \(0\le\gamma\le pL/2\). With \(t=\gamma/p\), the allowed arc has
the lift \([t-L,t]\), where \(0\le t\le L/2\).
Consecutive lifts in a surviving orbit satisfy \(z'=Bz+d\) with
\[
 d\in\mathbb Z\cap[-(B-1)t-L,\ BL-(B-1)t].
\]
Except at \(B=3,\ t=1/4\), this gives \(d\in\{0,1\}\).
A positive lift then grows until it leaves the allowed arc.
If \(d=1\), then \(z'\ge Bt-L\). Hence for \(t>L/B\) only \(0\)
survives, and at \(t=L/B\) any carry \(d=1\) in a surviving orbit
lands at \(0\). The latter survivor is contained in the preimages of
\(0\), so it is countable. In the exceptional case \(B=3,\ t=1/4\),
the nonzero carries are exactly the transitions between the two
endpoints \(-1/4,1/4\). Any surviving point is \(0\) or eventually
enters this two-cycle, since a trajectory with all carries zero
cannot stay bounded from a nonzero lift. This survivor is countable
as well.

If \(t<L/B\), then \(L-t>1/B\). Choose \(m\) large enough that
\(B^{-1}/(1-B^{-m})<L-t\). Concatenations of the two base-\(B\)
words \(0^m\) and \(10^{m-1}\) have every tail in
\([0,B^{-1}/(1-B^{-m})]\). Their negatives stay in the allowed arc,
and \(T_B^m\) on the aligned concatenations is conjugate to the full
shift on two symbols. Thus the entropy is positive. Undoing the
reflection and translation gives \([1/3,2/3]\bmod 1\) for tripling
and \([1/8,3/8]\bmod 1/2\) for quadrupling. The closed intervals
include the terminating-carry cases above.

\subsection{Doubling}

If \(p=2\), the restriction of \(T_2\) to the closed allowed
semicircle has only one possible image with two allowed preimages,
the common image of the endpoints. For each \(j\), the set of
depth-\(j\) preimages of either specified endpoint that stay in \(J\)
through that time has at most \(j+1\) elements. Indeed each
backward step adds at most one element, because at most one target
at the preceding depth can have two allowed preimages. These
preimages for \(0\le j<n\) divide the set surviving \(n\) tests into
\(O(n^2)\) intervals and cut points. To see this, on each complementary
interval the first \(n\) membership tests cannot change from all
allowed to a failed test without meeting a first boundary preimage.
On a surviving interval all of its first \(n\) images are arcs in
\(J\) and have lengths \(2^j\) times its own length. Its length is
therefore at most \(2^{-n}\). For any fixed \(\delta>0\) it can
contain only \(O_\delta(1)\) points of an \((n,\delta)\)-separated
set. The resulting polynomial bound gives entropy zero at every
\(\gamma\).

Now \(p\ge3\), so \(w\le1/3\). A hole whose closure contains \(0\)
is contained in one of width \(1/3\) whose closure also contains
\(0\). After reflection the latter hole has complement
\([v,v+2/3]\) in \([0,1]\) for some \(0\le v\le1/6\).
If \(v<1/6\), concatenate the words \(A=01\) and
\(D=001(01)^M\). No tail exceeds \(0.\overline{10}=2/3\). A tail
beginning in \(00\) must begin with \(D\), whose value tends to
\(0.00\overline{10}=1/6\) as \(M\to\infty\). All other tails are
at least \(1/4\). For sufficiently large \(M\) the two distinct
equal-length blocks \(AD,DA\) therefore give a full shift for
an iterate of doubling in the survivor. At \(v=1/6\), use aligned
concatenations of \(01\) and \(10\). Their aligned tails lie in
\([1/3,2/3]\) and their other tails in \([1/6,5/6]\). These
constructions give positive entropy also for holes abutting \(0\).

For the remaining holes write \(I=(a,c)\) with
\(0<a<c<1,\ c-a=w\). If \(c<1/2\), the two equally long binary
words \(1^m\) and \(01^{m-1}\), for \(m\) sufficiently large, give
a full shift for \(T_2^m\) whose tails all exceed \(c\).
If \(c=1/2\) and \(a>1/4\), use \(0^m\) and \(10^{m-1}\).
A tail starting in \(0\) has value at most
\(1/[4(1-2^{-m})]\), which is less than \(a\) for sufficiently large
\(m\), and every tail starting in \(1\) is at least \(1/2\).
In contrast, the survivor of \((1/4,1/2)\) is countable. In a
terminating-convention binary expansion any occurrence of \(01\)
must have a suffix of all zeros to avoid the hole. An expansion
with no occurrence of \(01\) is eventually constant. Thus a
hole ending at \(1/2\) with \(a\le1/4\) has zero entropy.
Reflection gives the corresponding results on the other side
of \(1/2\).

For a hole strictly crossing \(1/2\), Glendinning and Sidorov
characterize positive dimension by \(c<\chi(a)\)
\cite[Theorems~1.2 and~2.13]{app-glendinning-sidorov}. The interval
and circle survivors have the same dimension after identifying
\(0\) and \(1\). Holes with \(a\le1/4\) or \(c\ge3/4\) contain \((1/4,1/2)\) or its reflection
and have zero entropy. When \(1/4<a<1/2<c<3/4\), the bound
\(\chi(a)-a\le1/4\) in \cite[Proposition~3.3]{app-glendinning-sidorov}
gives zero entropy for all the crossing holes at \(p=3,4\).
The preceding one-sided analysis includes the two abutting cases,
so the full condition at these two widths is
\(a\in[1/2-w,1/2]\), or
\(\gamma\in[p/2-1,p/2]\bmod p\).

At \(p=5\), holes abutting \(1/2\) have positive entropy.
For strictly crossing holes, reflect if necessary so that
\(3/10<a\le2/5\). We use the following consequences of the
plateau description of \(\chi\)
\cite[Theorem~2.13 and Proposition~3.3]{app-glendinning-sidorov}.
The first \(1/2\) renormalization interval is
\([1/3,5/12]\). Outside that interval, either \(\chi(a)-a=1/4\)
or the first renormalization denominator \(q_1\) is at least \(3\).
In the latter case the lower bound from the plateau and gap
formulas is
\[
 \chi(a)-a\ \ge\
 \frac14\prod_{j\ge1}(1-2^{-3\cdot 2^{j-1}})
 \ \ge\ \frac14\left(1-\sum_{j\ge1}2^{-3j}\right)
 =\frac{3}{14}>\frac15 .
\]
Here the factors in those formulas have exponents \(q_1\cdots q_j\),
with \(q_i\ge2\) for \(i>1\). On a plateau the smallest width is
at its right endpoint and equals the product through that level
times \(1/4\), as in the proof of Proposition~3.3. The gap values
are finite products of the same factors, and the limiting
values use the infinite product.

The words \(01,10\) give
\(\chi(a)=7/12\) on \([1/3,19/48]\), whose endpoints have binary
expansions \(0.\overline{01}\) and \(0.0110\overline{01}\).
Thus this plateau gives zero entropy exactly for
\(a\ge23/60\). The next \(1/2\) substitution uses \(0110,1001\)
and gives
\(\chi(2/5)=0.1001\overline{0110}=47/80\).
Monotonicity then gives \(\chi(a)\le47/80<19/48+1/5\) for
\(19/48\le a\le2/5\), so this entire subinterval has zero entropy.
Reflection gives the full range \(23/60\le a\le5/12\).
The strict positivity inequality at the boundary gives the
closed range \([23/12,25/12]\bmod 5\) in \(\gamma\).

For \(p\ge6\), the universal positive-dimension width of
\cite[Corollary~3.9]{app-glendinning-sidorov} is strictly greater
than \(1/6\). Indeed
\[
 \frac14\prod_{j\ge1}(1-2^{-2^j})
 \ \ge\ \frac14\left(1-\frac14-\frac1{16}
                         -\sum_{m\ge8}2^{-m}\right)
 =\frac{87}{512}>\frac16 .
\]
This handles every hole with \(0<a<c<1\) in that range, and the
wrapping argument covers the rest. This proves Table~\ref{tab:zones}.

\subsection{The seven numeral radices}

For \(Q=1\) and \(B=2\), we have \(g=2\) and \(u\mid p\).
The rows \(p=2,3,4,5\) therefore give
\(\{2,4\},\{2,8\},\{2,4,16\},\{2,32\}\), respectively.
For \(Q=1,\ B=3,\ p=2\) the same divisibility gives \(3,9\).
For \(P=3,\ Q=2,\ B=4\) we have \(g=2\) and \(u\mid3\), giving
\(2,8\). Their union is \(\{2,3,4,8,9,16,32\}\). The argument
classifies the entropy of the survivor for admissible branches.
Actual prefix matches for a specified sequence depend on its
particular branch orbits.

\section{Combinatorial applications}\label{sec:applications}

\subsection{Fixed shifts of the central binomial coefficient}

Let \(s,t\in\mathbb Z\) be fixed. For any fixed integer \(a\), expansion
of a finite product or its reciprocal gives
\[
 \frac{(N+a)!}{N!}
 =N^a\left(1+\frac{a(a+1)}{2N}+O_a(N^{-2})\right).
\]
Applying this to the three factorials and using Stirling's formula
for \(\binom{2n}{n}\) gives
\begin{equation}\label{eq:app-shifted-binomial}
 \binom{2n+s}{n+t}
 =\frac{2^s4^n}{\sqrt{\pi n}}
  \left(1-\frac{\gamma}{n}+O_{s,t}(n^{-2})\right),
 \qquad
 \gamma=\frac18+\frac{(s-2t)^2+s}{4}.
\end{equation}
For the unshifted expansion with error bounds, see \cite{app-brent}.
The numerator \((s-2t)^2+s\) is even. Thus
\(\gamma=1/8\bmod 1/2\), the endpoint of the quadrupling zone.
For the binary numeral radix, \(p=3/2\), \(e=3\), \(B=4\), and
Theorem~\ref{thm:transfer} gives the growing hit counts in
\([X,X^{3/2+\varepsilon}]\) asserted in Theorem~\ref{thm:headline}. Fixed negative shifts are allowed since
all the factorial arguments are nonnegative for sufficiently large
\(n\).

\subsection{Lattice walks and the central trinomial square}

\begin{lemma}\label{lem:app-walk-expansion}
Let \(S\) be chosen uniformly from a finite list of \(m\) steps in a
full-rank lattice \(\Lambda\subset\mathbb R^d\), with repeated steps
counted as distinct labels. Suppose \(\mathbb E S=0\) and the
differences of steps generate \(\Lambda\). For each fixed
\(a\in\Lambda\), the number \(U_n(a)\) of length-\(n\) walks from
\(0\) to \(a\) satisfies
\[
 U_n(a)=
 \frac{\operatorname{covol}(\Lambda)\,m^n}
      {(2\pi n)^{d/2}\sqrt{\det\Sigma}}\,
 \left(1-\frac{\gamma}{n}+O(n^{-2})\right),
 \qquad \Sigma=\mathbb E[SS^\top],
\]
with \(\gamma\in\mathbb Q\). For a symmetric step law and \(a=0\),
\begin{equation}\label{eq:app-symmetric-correction}
 \gamma=\frac d4-\frac18\operatorname{Var}(S^\top\Sigma^{-1}S).
\end{equation}
\end{lemma}
\begin{proof}
Work first in a lattice basis, so \(\Lambda=\mathbb Z^d\). The
generation hypothesis makes \(\Sigma\) positive definite. It also
implies \(|\phi(t)|<1\) off \(0\) on the Fourier torus, where
\(\phi(t)=\mathbb E e^{it^\top S}\). In fact equality in the triangle
inequality would make \(e^{it^\top S}\) the same on every step, forcing
\(t\) to annihilate all of \(\Lambda\). Fourier inversion gives
\[
 \frac{U_n(a)}{m^n}
 =\frac1{(2\pi)^d}\int_{[-\pi,\pi]^d}e^{-it^\top a}\phi(t)^n\,dt .
\]
Away from a fixed neighborhood of \(0\) the integral is exponentially
small. In that neighborhood \(|\phi(t)|\le e^{-c|t|^2}\) for some
\(c>0\). Write \(t=y/\sqrt n\), let \(M_j(y)=\mathbb E(y^\top S)^j\),
and put \(A(y)=y^\top a+M_3(y)/6\). On \(|y|\le n^{1/24}\), Taylor
expansion of the local logarithm of \(\phi\) through degree five and
then of the exponential gives
\begin{align*}
 e^{-it^\top a}\phi(t)^n
 &=e^{-M_2(y)/2}\biggl[1-\frac{iA(y)}{\sqrt n}
 +\frac1n\left(\frac{M_4(y)-3M_2(y)^2}{24}
                -\frac{A(y)^2}{2}\right)\\
 &\hspace{24mm}
 +\frac{P_3(y)}{n^{3/2}}
 +O\!\left(n^{-2}(1+|y|^{M})\right)\biggr].
\end{align*}
Here \(P_3\) is an odd polynomial and \(M\) is fixed. The logarithmic
remainder after scaling is \(O(n^{-2}|y|^6)\), and the remaining
terms in the exponent beyond the quadratic term tend uniformly to
zero on this ball. The exponential remainder is therefore bounded
by \(n^{-2}\) times a polynomial after the Gaussian is factored out.
The rest of the neighborhood contributes a Gaussian tail. The odd
terms integrate to zero. For \(Y\) a centered Gaussian of covariance
\(\Sigma^{-1}\), the relative coefficient is
\begin{equation}\label{eq:app-walk-correction}
 -\gamma=\mathbb E_Y\left[
     \frac{M_4(Y)-3M_2(Y)^2}{24}
       -\frac12\left(Y^\top a+\frac{M_3(Y)}6\right)^2
                      \right].
\end{equation}
In lattice coordinates the polynomial and the covariance matrix
have rational coefficients. The Gaussian moment formula then gives
\(\gamma\in\mathbb Q\), including for nonsymmetric step laws.
Changing from a lattice basis to Euclidean coordinates multiplies the
determinant normalization by \(\operatorname{covol}(\Lambda)\).

If the law is symmetric and \(a=0\), then \(A=0\).
Writing \(R=S^\top\Sigma^{-1}S\), the Gaussian fourth moment gives
\(\mathbb E_Y M_4(Y)=3\mathbb E R^2\), and
\(\mathbb E_Y M_2(Y)^2=d(d+2)\). Since \(\mathbb ER=d\),
\eqref{eq:app-walk-correction} gives \eqref{eq:app-symmetric-correction}.
The quadratic form defining \(R\) is invariant under the change
of basis.
\end{proof}

For \(d=2\) and \(m=2^h\), the lemma gives \(p=2\) and a rational
\(\gamma\) for each fixed endpoint. Every such correction belongs to
the doubling zone. In base two, \(e=2\) and Theorem~\ref{thm:transfer}
gives the claimed \(2+\varepsilon\) window. The argument uses the
third moments and the fixed displacement in
\eqref{eq:app-walk-correction}, so it applies to directed step
lists with zero mean. For the eight king moves in \(\mathbb Z^2\),
the mean is zero by symmetry and differences of steps include both
coordinate unit vectors.

For the central trinomial \(T_n=[x^n](1+x+x^2)^n\), take one
coordinate with step law uniform on \(\{-1,0,1\}\). Its variance
is \(2/3\) and \(R=3S^2/2\) has variance \(1/2\). By
\eqref{eq:app-symmetric-correction} the one-coordinate correction is
\(\gamma=3/16\), in agreement with the trinomial expansion in
\cite{app-wagner}. The square satisfies
\begin{equation}\label{eq:app-trinomial}
 T_n^2=\frac{3\,9^n}{4\pi n}
       \left(1-\frac{3}{8n}+O(n^{-2})\right).
\end{equation}
In base three we have \(p=e=2,\ B=3\), and
\(3/8\in[1/3,2/3]\). Thus the ternary hit count tends to infinity
in every window \([X,X^{2+\varepsilon}]\).

\subsection{The binary Franel cutoff}

For fixed \(r\ge2\) let \(F_r(n)=\sum_{j=0}^n\binom nj^r\), the
Franel family studied in \cite{app-straub-zudilin}. The expansion
below is a special case of the binomial-sum expansions of
McIntosh~\cite{app-mcintosh}. We give the calculation of its first
correction since it determines the hole.
\begin{equation}\label{eq:app-franel}
 F_r(n)=
 \frac{(2/\pi)^{(r-1)/2}}{\sqrt r}\,
 2^{rn} n^{-(r-1)/2}
 \left(1-\frac{(r-1)^2}{4rn}+O_r(n^{-2})\right).
\end{equation}
To obtain the correction, put \(y=(j-n/2)/\sqrt n\).
On \(|y|\le\sqrt{K\log n}\), Stirling's formula gives uniformly
\begin{align*}
 \binom nj
 &=2^n\sqrt{\frac{2}{\pi n}}\,e^{-2y^2}
 \left(1+\frac{-1/4+2y^2-(4/3)y^4}{n}
          +O_K\!\left(\frac{1+|y|^{12}}{n^2}\right)\right).
\end{align*}
The entropy term contributes \(-4y^4/(3n)\), the square-root
factor \(2y^2/n\), and the Stirling correction \(-1/(4n)\).
Off this range, the centered binomial Chernoff bound
\(2^{-n}\binom nj\le e^{-2y^2}\) follows from
\(\cosh(t/2)\le e^{t^2/8}\).
Taking \(K\) sufficiently large makes the omitted tail smaller
than the stated error after raising to the \(r\)-th power.
For \(f\) a polynomial times \(e^{-2ry^2}\), Poisson summation on
the mesh \(1/\sqrt n\) gives
\[
 \sum_{j\in\mathbb Z} f((j-n/2)/\sqrt n)
   =\sqrt n\left(\int_{\mathbb R} f(y)\,dy+O_f(e^{-c_f n})\right)
\]
for some \(c_f>0\). The bound is uniform in the mesh translation
because the nonzero Fourier terms are bounded by a Gaussian
times a polynomial. The moments of the normalized
Gaussian weight are \(1/(4r)\) and \(3/(16r^2)\) in degrees two
and four, respectively. The relative coefficient is therefore
\[
 r\left(-\frac14+\frac{2}{4r}-\frac43\frac{3}{16r^2}\right)
 =-\frac{(r-1)^2}{4r}.
\]
The leading Gaussian integral is \(\sqrt{\pi/(2r)}\), giving
\eqref{eq:app-franel}.

Here \(p=(r+1)/2\) and \(\gamma=(r-1)^2/(4r)\).
When \(r=2\), Vandermonde's identity gives the unshifted central
binomial coefficient. For \(r=3,5,7\) the respective pairs
\((p,\gamma)\) are \((2,1/3),(3,4/5),(4,9/7)\), all in the
doubling zones. They give \(e=p\) and window exponents
\(2+\varepsilon,3/2+\varepsilon,4/3+\varepsilon\), respectively.
For even \(r\ge4\), the branch map is quadrupling and \(p\ge5/2\),
so \eqref{eq:app-inverse-tree} gives positive entropy. At
\(r=9\) we have \(p=5,\ \gamma=16/9<23/12\), and for odd \(r\ge11\)
we have \(p\ge6\). Both cases have positive entropy by
Table~\ref{tab:zones}. This gives the exact entropy cutoff for
the unpowered Franel family in base two.

\subsection{Perfect-power numeral radices}

Let \(W_n\) count \(n\)-step returns in the following
eight-dimensional walk. Partition the coordinate labels into a
triangle and a pentagon. For each of the eight graph edges
\(\{i,j\}\), allow the four distinct moves \(\pm e_i\pm e_j\)
with independent signs. The step lattice consists of integer
vectors with even coordinate sum on each block. Indeed the edge
differences \(e_i-e_j\) generate the zero-sum lattice on each
connected block, and combining two signs on an incident edge
gives \(2e_i\) for any vertex. The block parity conditions are
independent, so the lattice has index \(4\) in \(\mathbb Z^8\).
To check strong aperiodicity, quotient this lattice by the subgroup
generated by pairwise differences of steps. All steps have the
same class, which generates the quotient. A step and its negative
give a two-step return, and the three oriented differences around
the triangle give a three-step return. These make both twice and
three times the common class zero, so the quotient is trivial.

The law is symmetric, and \(\Sigma=I_8/4\) since each coordinate is
changed on two of the eight edges and the two signs are
independent. Thus \(R=8\) identically and
Lemma~\ref{lem:app-walk-expansion} gives
\begin{equation}\label{eq:app-32walk}
 W_n=\frac{64\,32^n}{\pi^4 n^4}
       \left(1-\frac2n+O(n^{-2})\right).
\end{equation}
In particular the base-\(32\) parameters are \(h=1,p=5,\gamma=2\).
The other square expansions needed below are
\begin{align*}
 \binom{2n}{n}^2&=\frac{16^n}{\pi n}
                  \left(1-\frac1{4n}+O(n^{-2})\right),\\
 F_3(n)^2&=\frac{4\,64^n}{3\pi^2 n^2}
                  \left(1-\frac2{3n}+O(n^{-2})\right),\\
 F_4(n)^2&=\frac{2\,256^n}{\pi^3 n^3}
                  \left(1-\frac9{8n}+O(n^{-2})\right).
\end{align*}
The last square has \(p=4,\gamma=9/8\), and \((B,e)\) equals
\((2,4),(2,2),(2,1)\) in numeral bases \(2,4,16\), respectively.
Its exact success arc is \((9/32,17/32)\), with
\(\gamma\) in the zero-entropy zone for \(p=4\). Theorem~\ref{thm:transfer}
gives counts tending to infinity in windows of exponents
\(4/3+\varepsilon\) in binary and \(2+\varepsilon\) in base four.

In each of the first five rows of Table~\ref{tab:app-inner}, the
scale ratio \(b^{1/p}\) is the integer \(B\), so \(e=1\).
The indicated hole has rational endpoints and closure strictly
inside the exact success arc from the corresponding expansion.
Each associated \(C^{1/p}\) is irrational since \(C\) is a nonzero
rational multiple of a negative integral power of \(\pi\)
\cite{app-lindemann}. Hence the single-branch clause of
Theorem~\ref{thm:transfer} gives infinitude in all five rows. The
binary and ternary applications already proved supply the
realizations in the two smallest radices.

\begin{table}[ht]
\centering
\small
\begin{tabular}{@{}rcllcl@{}}
\toprule
\(b\) & \(B\) & \(U_n\) & \((l,r,D)\) & \(s\) & Cyclic components \\
\midrule
4  & 2 & \(\binom{2n}{n}^2\) & \((3,9,16)\)       & 1 & \((0), (15)\) \\
8  & 2 & \(F_3(n)^2\)        & \((8,17,32)\)      & 1 & \((0), (31)\) \\
9  & 3 & \(T_n^2\)           & \((8,27,40)\)      & 1 & \((0), (39)\) \\
16 & 2 & \(F_4(n)^2\)        & \((145,271,512)\)  & 2 & \((0), (34,68,136,273), (511)\) \\
32 & 2 & \(W_n\)             & \((13,19,32)\)     & 3 & \((0), (6,12,25,19), (10,21), (31)\) \\
4  & 2 & \(F_4(n)^2\)        & \((9,17,32)\)      & 2 & \((0), (2,4,8,17), (31)\) \\
\bottomrule
\end{tabular}
\caption{Cyclic components for the grid hole \((l/D,r/D)\), listed
in cycle order. The last row uses the exact success arc. The
integer \(s\) is the largest number of distinct cyclic components
on a directed path.}
\label{tab:app-inner}
\end{table}

The table can be reproduced by retaining cells
\(i\in\{0,\ldots,l-1\}\cup\{r,\ldots,D-1\}\) and forming the
edge list \(i\to(Bi+t)\bmod D\) over \(0\le t<B\), whenever the
target is retained. For each retained \(i\), start with \(R_i=\{i\}\)
and repeatedly add targets of outgoing edges from \(R_i\) until
it stabilizes. The component of \(i\) consists of \(j\in R_i\)
with \(i\in R_j\). Counting internal labelled edges gives one
per vertex in each listed component. For \(s\), order distinct
cyclic components by reachability between representatives and take
the longest chain. In the first three rows the lower loop reaches
only lower cells, and the only retained predecessor of cell \(0\)
is \(0\) itself. Hence neither loop is reachable from the other.
In base nine, the alternate inner hole \((10/48,32/48)\)
also has simple cycles but allows both on one directed path,
so we use the displayed hole for the window estimate.

\subsection{Explicit single-branch window exponents}

If \(\beta=c/\alpha^d\) with \(c>0\) rational and \(d\) a
positive integer, then
\(\mu(\alpha)\le d\,\mu(\beta)\) for irrational \(\alpha,\beta>0\).
Indeed any rational approximant \(a/q\) sufficiently close to
\(\alpha\) gives an approximant \(c(q/a)^d\) with denominator
\(O(q^d)\), and the map \(x\mapsto c/x^d\) is locally Lipschitz
at \(\alpha\). The irrationality-exponent inequality follows by
applying this to infinitely many approximants. For the base-four
binomial square and the base-nine trinomial square,
\(\alpha^2\) is a rational multiple of \(1/\pi\). For the
base-eight Franel square, \(\alpha^3=4/(3\pi^2)\), so one can
take \(\beta=\zeta(2)=\pi^2/6\) with \(d=3\). In all three rows
\(s=1\). The bounds \(\mu(\pi)\le7.103205334137\ldots\)
\cite{app-zeilberger-zudilin} and
\(\mu(\zeta(2))\le5.095412\) \cite{app-zudilin}, together with
Theorem~\ref{thm:transfer}, give at least one hit in every
sufficiently large \([X,X^\rho]\) for
\begin{equation}\label{eq:app-single-branch-exponents}
 \begin{cases}
  \rho>14.206411,& \binom{2n}{n}^2\text{ in base }4
                   \text{ and }T_n^2\text{ in base }9,\\
  \rho>15.286236,& F_3(n)^2\text{ in base }8 .
 \end{cases}
\end{equation}
Taking disjoint power windows gives \(\gg\log\log X\) hits up to
\(X\) for each of these three examples.

\subsection{A lower bound for the number of prefixes}\label{sec:app-counting}

\begin{theorem}\label{thm:app-counting}
Suppose \(U\) satisfies \eqref{eq:aligned}, with \(e>1\) in
\eqref{eq:branch}. Assume that \(I\) has rational endpoints and
\(K_B(I)\) has zero entropy, or that there is an open arc \(J\) with
rational endpoints, \(\overline J\subset I\), and \(K_B(J)\) of zero
entropy. Then for all sufficiently large \(X\),
\begin{equation}\label{eq:app-counting}
 H_b(U;1,X)\gg(\log\log X)(\log\log\log X).
\end{equation}
In particular this holds in base two for \(B_n\) and its fixed shifts
\(\binom{2n+s}{n+t}\), and for \(F_3,F_5,F_7\). It holds in base
three for \(T_n^2\) and in bases two and four for \(F_4^2\).
It also holds in base two for the length-\(n\) walk count \(U_n(z)\)
from \(0\) to any fixed \(z\in\Lambda\) for a centered list of
\(2^h\) steps in a rank-two lattice \(\Lambda\), with integer \(h\ge1\),
whose pairwise differences generate \(\Lambda\). Repeated steps
count as distinct choices.
\end{theorem}

The finite graph in Section~\ref{sec:transfer} already bounds the
number of transitions by a constant times one plus the number of
successful tests. We quantify Lemma~\ref{lem:transitions} on
geometrically spaced windows using the interval theorem of Evertse
and Ferretti.

\begin{lemma}\label{lem:app-count-transitions}
Let \(B,d,v\) be as in Lemma~\ref{lem:transitions}. Fix \(E\ge0\),
a compact set \(\mathcal C\subset\mathbb R\setminus\{0\}\),
\(\eta>d/(d-1)\), and an integer \(A>\eta\). Put \(u_j=A^j\),
\(w_j=\lfloor\eta u_j\rfloor\) for \(j\ge0\). In each window choose
\(c_j\in\mathcal C\) and integer vectors \(W_m\) satisfying
\(\|W_m-c_j B^m v\|_\infty\le E\) for \(u_j\le m\le w_j\).
If \(s_j\) is the number of edges \(u_j<m\le w_j\) with
\(W_m\ne BW_{m-1}\), then there is a constant \(C_0>0\), depending
only on the fixed data, such that
\begin{equation}\label{eq:app-count-exceptions}
 \#\{j:s_j\le k\}\le\exp(C_0(k+1))\qquad(k\ge0,\ k\in\mathbb Z).
\end{equation}
\end{lemma}

\begin{proof}
Write \(u=u_j,\ w=w_j,\ c=c_j\), and let \(a_1<\cdots<a_s\) be
the transition times in this window. Set \(h_i=W_{a_i}-BW_{a_i-1}\).
These vectors belong to a fixed finite set since
\(\|h_i\|_\infty\le(B+1)E\). Use the vector
\(Z=(Y_1,\ldots,Y_d,P_1,\ldots,P_s)\) in dimension \(N=d+s\), with
\(Y=B^{w-u}W_u\) and \(P_i=B^{w-a_i}\). Use the same forms as in
Lemma~\ref{lem:transitions}, now with one power for each transition.
At infinity these are \(F_1,\ F_r-(v_r/v_1)F_1\) for \(2\le r\le d\),
and \(P_i\) for \(1\le i\le s\), where \(F=Y+\sum_i h_iP_i\).
At finite places take the coordinate forms. All local determinants
are \(1\). Telescoping gives \(F(Z)=W_w\).

Choose \(1/d<\lambda<1-1/\eta\), put \(\zeta=d\lambda-1\) and
\(g=\zeta/(4N)\). For large \(u\) we have \((w-u)/w\ge\lambda\).
Round each \((w-a_i)/w\) down to a multiple \(l_i\) of \(g\).
For \(\ell\mid B\) write
\(\beta_\ell=\nu_\ell(B)\log\ell/\log B\), so that
\(\sum_{\ell\mid B}\beta_\ell=1\). The bounds for the evaluations
at \(Z\) have the following exponents \(e_{i,p}\) to base \(B^w\),
with a uniform constant at infinity and unit constant at finite places.
\[
 \begin{array}{c|cc}
 p & \text{first \(d\) forms} & \text{form for \(P_i\)}\\ \hline
 \infty & 1,0,\ldots,0 & l_i+g\\
 \ell\mid B & -\beta_\ell\lambda\ \text{each} & -\beta_\ell l_i\\
 \ell\nmid B & 0\ \text{each} & 0
 \end{array}
\]
Indeed \(F_1(Z)=O(B^w)\), the other \(F\)-forms on \(Z\) are
\(O(1)\), and \(Y\) is divisible by \(B^{w-u}\). The sum of all
entries is \(S=-\zeta+sg\le-3\zeta/4\).


Take a fixed Galois number field \(K\subset\mathbb C\) containing
the \(v_r\). For each infinite place \(\nu\), take an embedding
\(\sigma_\nu\) representing it and apply \(\sigma_\nu^{-1}\) to the
real form coefficients. The choice between conjugate embeddings is
immaterial since these coefficients are real. At a place \(\nu\)
above \(p\), put \(d_\nu=[K_\nu:\mathbb Q_p]/[K:\mathbb Q]\), with
\(\mathbb Q_\infty=\mathbb R\). Since \(Z\) is rational, the table
applies to the normalized absolute values \(\|\cdot\|_\nu\) with
exponents \(d_\nu e_{i,p}\). Let \(\bar e_p=N^{-1}\sum_i e_{i,p}\)
and set
\[
 \kappa_{i,\nu}=\tfrac14 d_\nu(e_{i,p}-\bar e_p),\qquad Q=B^{4w}.
\]
The weights have finite support and zero sum at each place.
Their local maxima sum to at most \((2+g)/4<1\), since
\(\sum_{\nu\mid p}d_\nu=1\), the range at infinity is at most
\(1+g\), and the range at \(\ell\mid B\) is at most \(\beta_\ell\).
In the notation of \cite[(2.4)--(2.15)]{cnt-evertse-ferretti},
\(\Delta_{\mathcal L}=1\), and the twisted height satisfies
\[
 H_{\mathcal L,\kappa,Q}(Z)
 =\prod_\nu\max_i\bigl(\|L_i^{(\nu)}(Z)\|_\nu Q^{-\kappa_{i,\nu}}\bigr)
 \ll B^{wS/N}\le B^{-3\zeta w/(4N)}.
\]
In particular it is at most \(Q^{-\delta}\) for
\(\delta=\zeta/(8N)\in(0,1)\) and \(w\) at least a sufficiently
large constant times \(N\).

There are at most \(R=([K:\mathbb Q]+1)N\) distinct forms.
Their coefficients have a fixed common integral denominator
and bounded conjugates, uniformly in \(s\) and the \(h_i\).
Hadamard's inequality at the infinite places bounds the
determinants in \(H_{\mathcal L}\) by \((C_1\sqrt N)^N\)
before local normalization. At finite places clearing the fixed
denominator bounds their product by its \(N\)-th power.
Thus \(\log H_{\mathcal L}=O(N\log(2N))\), and the same
coefficient bounds give \(\max_{\nu,i} H_2(L_i^{(\nu)})\le C_2\sqrt N\)
for the Euclidean height of \cite[\S6]{cnt-evertse-ferretti}.

Let \(T\) be the distinguished proper subspace in
\cite[(2.19)--(2.21)]{cnt-evertse-ferretti}. For a subspace \(U\),
its weight \(w(U)\) sums the minimum weight of a basis of the
restricted forms at each place. The Galois action permutes the
places and their systems of forms and weights, by the inverse
embedding construction at infinity. Since \(T\) is defined over
\(K\) and uniquely maximizes \(w(U)/(N-\dim U)\) among subspaces
of minimal dimension attaining the maximum, it is defined over
\(\mathbb Q\).

Write \(Y_0=\{(Y,0)\}\). If \(U\supseteq Y_0\) is proper and has
dimension \(r\), select \(r-d\) of the \(P_i\) giving a coordinate
basis of the dual of the projection of \(U\) onto the \(P\)
coordinates. Together with the first \(d\) forms, these give a
basis of restricted forms at each place, because the first \(d\)
restrict to a basis on \(Y_0\). The terms in \(l_i\) for the
selected \(P_i\) cancel across places. The resulting total raw
weight is \(-\zeta+(r-d)g=S-(N-r)g\). After centering,
\[
 w(U)\le \tfrac14\bigl(S(1-r/N)-(N-r)g\bigr)<0.
\]
Since \(w(0)=0\), we have \(T\not\supseteq Y_0\).

By \cite[Proposition~17.5]{cnt-evertse-ferretti},
\(H_2(T)\le (C_2\sqrt N)^{4^N}=:D_N\), increasing \(C_2\) to at
least \(1\) if necessary. There is an integer relation
\(a\cdot Y+b'\cdot P=0\) on \(T\) with \(a\ne0\) and all
coefficients of magnitude at most \(D_N\). Indeed, a primitive
integer Pl\"ucker vector of \(T\) has Euclidean norm \(H_2(T)\).
Take a coordinate of its wedge with \((Y,P)\) which is nonzero
on some coordinate vector in \(Y_0\). For \(T=0\) a \(Y\)-coordinate
itself suffices. By rational independence \(a\cdot v\ne0\).
After clearing the fixed denominators of the \(v_i\), its
algebraic norm is a nonzero integer with all other conjugate
factors \(O(D_N)\). Hence
\[
 |a\cdot v|\gg D_N^{-([K:\mathbb Q]-1)}.
\]
If \(Z\in T\), dividing the relation by \(B^w\) gives
\[
 |c\,a\cdot v|\ll N D_N B^{-u},
\]
because \(Y/B^w=cv+O(B^{-u})\) and \(P_i/B^w=B^{-a_i}\).
The positive lower bound for \(|c|\) on \(\mathcal C\) shows that
\(Z\notin T\) for \(u>C_3\,4^N\log(2N)\), with a fixed \(C_3\).

For fixed \(s\), a fixed ordered list of the \(h_i\), and a fixed
profile \((l_i)\), the forms and weights are fixed.
Evertse and Ferretti's interval theorem
\cite[Theorem~2.3]{cnt-evertse-ferretti} puts all \(Q\) with a
point of twisted height at most \(Q^{-\delta}\) outside \(T\) in
\([1,\max(H_{\mathcal L}^{1/R},N^{1/\delta})]\) and at most
\[
 10^5 4^N N^{10}\delta^{-2}\log(3\delta^{-1}R)
\]
intervals of the form \([Q_*,Q_*^\omega]\), where
\(\omega=\delta^{-1}\log(3R)\). The logarithm of the cutoff for
the initial interval is \(O(N\log(2N))\). This cutoff, the
small-height threshold and the threshold for \(Z\notin T\) together
omit at most \(O(N)\) early starts, since \(u_j=A^j\).
Each other interval contains at most \(1+\log_A\omega\) of the
parameters \(B^{4w_j}\), since \(w_{j+1}\ge A w_j\). Thus at most
\(\exp(O(N))\) starts realize a fixed list and profile.

There are \(\exp(O(s))\) possible lists. Since the grid profile
is weakly decreasing and has entries in
\(\{0,g,\ldots,\lfloor1/g\rfloor g\}\), there are at most
\(\binom{s+\lfloor1/g\rfloor}{s}=\exp(O(N))\) profiles.
Summing over \(0\le s\le k\) proves \eqref{eq:app-count-exceptions}.
\end{proof}

\begin{proof}[Proof of Theorem~\ref{thm:app-counting}]
Choose \(\eta>e/(e-1)\) and an integer \(A>\eta\). In the windows
of Lemma~\ref{lem:app-count-transitions}, let \(H_j\) count all
successful candidate tests across the \(e\) branches. The construction
in Section~\ref{sec:transfer} gives integer approximants to
\(c B^m(1,\tau,\ldots,\tau^{e-1})\) with a fixed bounded error and at
most \(C_4(H_j+1)\) vector transitions in the window. Here \(c\)
is \(C^{1/p}\) times the grid denominator and the common multiple
clearing the periodic tails. These factors are fixed across
windows. The degree calculation preceding \eqref{eq:line} gives
the required rational independence. By \eqref{eq:app-count-exceptions},
the number of windows with \(H_j\le k\) is at most
\(\exp(C_5(k+1))\) for a fixed \(C_5>0\).

All candidate indices in the \(j\)-th window are at most
\(C_6 B^{w_j}\) for a fixed \(C_6>0\). After a fixed initial
segment, they are distinct across all times and branches, as the
continuous indices in \eqref{eq:line} grow in order by the fixed
factor \(\tau>1\) and the candidate differs from its continuous
index by \(O(1)\). Thus windows with \(C_6B^{w_j}\le X\) supply
\(J\asymp\log\log X\) disjoint collections of regular candidate
tests up to \(X\). The \(r\)-th smallest \(H_j\) among these windows
is at least \((\log r)/C_5-1\). Summing over \(1\le r\le J\)
gives \(\gg J\log J\) successful tests, and each is an exact
prefix hit. The identical summation for \eqref{eq:app-count-exceptions}
gives \(\gg(\log M)(\log\log M)\) transitions in completed windows
through branch time \(M\).

The expansions \eqref{eq:app-shifted-binomial},
\eqref{eq:app-trinomial}, and \eqref{eq:app-franel}, and
Lemma~\ref{lem:app-walk-expansion} verify the hypotheses for
the first four families. Their respective \((B,e)\) values
are \((4,3)\) for the shifted binomial coefficients, \((3,2)\)
for the trinomial square, \((2,2),(2,3),(2,4)\) for \(F_3,F_5,F_7\),
and \((2,2)\) for the planar walks. For \(F_4^2\) the parameters
computed above give \((B,e)=(2,4),(2,2)\) in bases two and four,
respectively, with its rational hole in the doubling zone.
\end{proof}

\begin{remark}
The order \(J\log J\) is sharp for an exponential exceptional-window
count alone. Indeed, given \(C_7\ge1\) and \(C_8>0\), the abstract
counts \(h_j=\max(0,\lceil(\log j-\log C_7)/C_8\rceil)\),
\(j\ge1\), satisfy \(\#\{j:h_j\le k\}\le C_7 e^{C_8 k}\)
and sum to \(\asymp J\log J\) through \(J\).
There is also a single event set in integer branch time with
this order for all fixed multiplicative window ratios. For large
integers \(m\), put an event at \(m\) whenever
\(\lfloor f(m)\rfloor>\lfloor f(m-1)\rfloor\), where
\(f(t)=(\log t)(\log\log t)\). Since \(0<f'(t)<1\) eventually,
the event count through \(M\) is \(f(M)+O(1)\).
For every fixed \(\eta>1\), the count in
\((u,\lfloor\eta u\rfloor]\) is
\((\log\eta)\log\log u+O_\eta(1)\), by applying the mean value
theorem to \(y\log y\) with \(y=\log u\).
On a schedule with consecutive starts in ratio at least
\(q>1\), a bound of \(k\) on the count therefore forces
\(\log\log u\ll_\eta k+1\) and permits at most
\(\exp(O_{\eta,q}(k+1))\) starts. The same upper bound on
exception counts holds if the window ratios vary with a fixed
lower bound greater than \(1\), by retaining an initial subwindow
of fixed ratio. The event count remains of order
\((\log M)(\log\log M)\), or \((\log\log X)(\log\log\log X)\)
on the scale \(M\asymp\log X\). This shows sharpness for the
counting information about windows.
\end{remark}

\section{Metric spectra on algebraic lines}\label{sec:metric}

Let \(d\geq2\) and let \(v=(v_1,\ldots,v_d)\) have positive real algebraic
coordinates linearly independent over \(\mathbb Q\). For \(x>0\), its
\(j\)-th radix-\(B\) digit after the point is
\(\lfloor B^jx\rfloor-B\lfloor B^{j-1}x\rfloor\). For \(w>u\), a word with
integer value \(b\) in positions \(u+1,\ldots,w\) is equivalent to
\[
 \{B^ux\}\in[bB^{u-w},(b+1)B^{u-w}).
\]
Thus a prescribed word in each expansion of \(cv_i\) determines a
half-open box of fractional parts. Throughout this section \(\rho>1\),
the end of a radix window starting at \(u\) is \(\lfloor\rho u\rfloor\),
and the level conditions concern infinitely many integer starts.
Dimension statements are local on every nonempty bounded open interval
of \(c>0\), with the empty set assigned dimension zero.

For two bases \(B_1,B_2\geq2\), put
\[
 Q=N=B_1^u,\qquad u_2=\lfloor\log_{B_2}Q\rfloor,\qquad M=B_2^{u_2},
 \qquad G=\gcd(N,M).
\]
The two starts are \(u\) and \(u_2\). In particular \(M\asymp Q\), so
both targets have widths \(\asymp Q^{1-\rho}\) in fractional-part
coordinates and \(\asymp Q^{-\rho}\) in \(cv_i\). Write
\begin{equation}\label{eq:met-exponents}
 \begin{split}
 \lambda
 &=\sum_{p\ \mathrm{prime}}\min\left\{
       \frac{\nu_p(B_1)\log p}{\log B_1},
       \frac{\nu_p(B_2)\log p}{\log B_2}
     \right\},\\
 D_f&=2d-(2d-1)\rho,\qquad
 D_g=d\lambda-(d-1)\rho.
 \end{split}
\end{equation}
Taking valuations of \(N\) and \(M\) gives
\(\log G=\lambda\log Q+O(1)\). The two lists of normalized valuations
each sum to \(1\), so \(0\leq\lambda\leq1\), with \(\lambda=1\) exactly
when the lists agree. Agreement at any prime in their common support
then gives \(\log B_1/\log B_2\in\mathbb Q\). Conversely, a common
positive power gives agreement at every prime. In particular, the
two-base cutoff
\[
 \rho_*=\max\left\{\frac{2d}{2d-1},\frac{d\lambda}{d-1}\right\}
\]
is strictly smaller than \(d/(d-1)\) for multiplicatively independent
bases. 
\subsection{Dual bounds and counting along a segment}

We apply the Subspace Theorem in the form used in the proof of
Lemma~\ref{lem:transitions}. The one-denominator estimate uses the
real place \cite{schmidt}; for two bases we also use the primes dividing
\(B_1B_2\) \cite{met-schlickewei}.

\begin{lemma}\label{lem:met-dual}
For every \(\eta>0\) and nonzero \(k\in\mathbb Z^d\),
\[
   |k\cdot v|\gg_{\eta,v}\|k\|_\infty^{-(d-1)-\eta}.
\]
For the two bases in \eqref{eq:met-exponents}, let \(t=\rho-1\).
If \(t<1-\lambda\), then for every fixed \(C,\eta>0\) and all
sufficiently large \(u\), every nonzero \((k,h)\in\mathbb Z^{2d}\)
with \(\|(k,h)\|_\infty\leq C Q^t\) satisfies
\begin{equation}\label{eq:met-two-dual}
  |(Nk+Mh)\cdot v|\ \geq\ Q^{D_f-\eta}.
\end{equation}
\end{lemma}

\begin{proof}
For the first assertion, suppose the bound with constant \(1\) fails
on infinitely many integer vectors. Use \(k\cdot v\) and all but one of
the coordinate forms at the real place. The product is
\(O(\|k\|_\infty^{-\eta})\), so the Subspace Theorem places infinitely
many of the vectors in a single proper rational subspace. On a
rational subspace of rank \(m\), use any \(m\) coordinate projections
giving an isomorphism onto \(\mathbb Q^m\). The new coordinates are
integers of comparable norm, and the restricted coefficients of the
form are rationally independent because \(v\) is. The product of the
restricted form and \(m-1\) coordinate forms still decays by a fixed
negative power. Repeating the descent gives a contradiction in rank
one, where a nonzero linear form on integer vectors has positive
minimum modulus. The finitely many remaining exceptions are absorbed
in the implicit constant.

For \eqref{eq:met-two-dual} let \(N'=N/G\) and \(M'=M/G\), which are
coprime and each \(\asymp Q^{1-\lambda}\). Suppose infinitely many
levels violate the estimate. During the descent the two blocks \(k,h\)
range in fixed rational subspaces of ranks \(m_1,m_2\), initially
both \(d\). Project each onto integral coordinates as above, denoted
by \(k',h'\), and set \(m=m_1+m_2\). Apply the Subspace Theorem to
\(z=(N'k',M'h')\), for \(m\geq2\). At infinity use the form evaluating
to \((N'k+M'h)\cdot v\), and all but one coordinate form chosen so
that the forms are independent. At the primes dividing \(B_1B_2\)
use the coordinate forms. The full product is bounded by a constant
times
\[
 Q^{1-\lambda-(2d-1)t-\eta}\,
 Q^{(m-1)(1-\lambda+t)}\,
 Q^{-m(1-\lambda)}
 = Q^{-(2d-m)t-\eta}.
\]
Here the finite-place bound uses the \(N'\) or \(M'\) dividing every
coordinate in its block. The height of \(z\) is at most a fixed power
of \(Q\), so some fixed nontrivial rational relation holds along
infinitely many of the levels. After clearing denominators it reads
\(N'A(k')+M'H(h')=0\), where \(A,H\) are integral linear forms.
Coprimality gives \(M'\mid A(k')\) and \(N'\mid H(h')\). Both values
are \(O(Q^t)\), hence both vanish for large \(u\) since \(t<1-\lambda\).
At least one block loses rank. In total rank one the nonzero form has
modulus \(\gg Q^{1-\lambda}\), which contradicts the first factor of
the displayed bound. This ends the descent.
\end{proof}

\begin{lemma}\label{lem:met-segment}
Let \(Q=Q_u\) grow at least geometrically with \(u\). Suppose
\(a=a_u\in\mathbb R^k\) has \(\|a\|_\infty\asymp Q^\beta\) with
\(0<\beta<\rho\). Set
\[
 \delta=Q^{\beta-\rho},
 \qquad D=k\beta-(k-1)\rho.
\]
Assume \(D>0\) and, for every fixed \(C,\eta>0\), that
\[
 |a\cdot z|\geq Q^{D-\eta}\qquad
 (0\ne z\in\mathbb Z^k,\ \|z\|_\infty\leq C/\delta)
\]
eventually in \(u\). For any compact parameter interval and any
translate of a box of side lengths \(O(\delta)\) modulo
\(\mathbb Z^k\), the condition that \(ca\) lands in the box can be
covered by \(O(Q^D)\) intervals of length \(O(Q^{-\rho})\), uniformly
in the translation. If one box of side lengths \(\asymp\delta\) is
specified at each of infinitely many levels, the limsup has local
dimension at least \(D/\rho\).
\end{lemma}

\begin{proof}
Fix \(b>0\) and \(0<\epsilon<D\), put \(L=Q^{-D+\epsilon}\), and
consider the convex body
\[
 K=[-L,L]a+[-b\delta,b\delta]^k.
\]
Its width in an integral direction \(z\ne0\) is
\(2L|a\cdot z|+2b\delta\|z\|_1\). If the second term is at most the
flatness constant in dimension \(k\), then \(\|z\|_\infty\ll1/\delta\)
and the first term tends uniformly to infinity by the hypothesis
with \(\eta<\epsilon\). Hence every translate of \(K\) meets
\(\mathbb Z^k\) by the flatness theorem \cite{met-flatness}.

For an upper count, take a bounded parameter interval \(I\). For
some translation vector \(\theta\) and fixed \(C_1\), the relevant
lattice points lie in
\[
 I a-\theta+[-C_1\delta,C_1\delta]^k.
\]
Since \(K+\mathbb Z^k=\mathbb R^k\), a measurable fundamental domain
can be chosen inside \(K\). Indeed, for each \(x\in[0,1)^k\) choose
the first integer translate \(x+z\) in \(K\) from the finite list of
possible \(z\). This selects one representative per integer coset
and preserves volume one. The translates of this domain by the
relevant lattice points are disjoint and are contained in the
displayed set plus \(K\). Its extent parallel to \(a\) is
\(O(Q^\beta)\) and its orthogonal projection is contained in a ball
of radius \(O(\delta)\). Its volume is therefore
\(O(Q^\beta\delta^{k-1})=O(Q^D)\). Each lattice point contributes
at most one interval, since the box is convex, and one coordinate
with \(|a_i|=\|a\|_\infty\) bounds the length by
\(O(\delta/\|a\|_\infty)=O(Q^{-\rho})\).

For a lower bound, use the midpoint \(\theta\) of the specified box
and choose \(b\) uniformly small compared with the side lengths
divided by \(\delta\). The translate \(ca-\theta+K\) contains an
integer point. Thus for some \(c'\) with \(|c'-c|\leq L\), the point
\(c'a\) is within \(b\delta\) of the midpoint modulo
\(\mathbb Z^k\). An interval centered at \(c'\) with radius
\(r=\gamma Q^{-\rho}\), for a fixed sufficiently small \(\gamma>0\),
still lands strictly inside the box. Choose finitely many such
centers per level by sampling a mesh of spacing at most \(L\) on a
closed subinterval of the parameter interval. For
\(0<s<D/\rho\) choose \(\epsilon\) so that \(D-\epsilon>\rho s\).
Then \(r^s/L\to\infty\), and the intervals enlarged to radius
\(r^s\) cover the sampled subinterval at every sufficiently large
selected level. The one-dimensional Mass Transference Principle
of Beresnevich and Velani \cite[Theorem~2]{met-mtp} gives dimension
at least \(s\) for the limsup of the original intervals inside it.
Let \(s\) tend to \(D/\rho\).
\end{proof}

For \(a=Hv\) and \(H\asymp Q^\beta\), the first assertion of
Lemma~\ref{lem:met-dual} gives the hypothesis of
Lemma~\ref{lem:met-segment} with \(k=d\) and
\(D=d\beta-(d-1)\rho>0\). Indeed for \(\|z\|_\infty\ll1/\delta\)
it gives
\[
 H|z\cdot v|\gg Q^{D-\eta'(\rho-\beta)}
\]
for every \(\eta'>0\). This also applies to boxes with a fixed
additive change in the window's end position because only the
constants in the side lengths change. For \(a=(Nv,Mv)\),
Lemma~\ref{lem:met-dual} supplies the same hypothesis with
\(k=2d\), \(\beta=1\), and \(D=D_f\) in the range
\(\rho<2-\lambda\) and \(D_f>0\).

\subsection{Proof of the dimension and strict cutoff assertions}

For one base take \(Q=B^u\) and \(H=Q\). When
\(D=d-(d-1)\rho>0\), Lemma~\ref{lem:met-segment} gives the lower
bound \(D/\rho\) for any choice of one word per selected level,
including the zero word. If at most \(\exp(o(u))\) words are
allowed at level \(u\), the \(s\)-power sum of the interval lengths
in the upper cover is bounded by a constant times
\(\sum_u Q^{D-\rho s+o(1)}\). For \(s>D/\rho\) this converges,
and covering the limsup by arbitrarily late tails gives
\(\mathcal H^s=0\). For \(\rho\geq d/(d-1)\), truncate each word
at \(\lfloor\sigma u\rfloor\) with \(1<\sigma<d/(d-1)\), keeping
at most the same number of words. Letting \(\sigma\) increase
to \(d/(d-1)\) gives the dimension-zero upper bound. Thus the
one-base dimension in Theorem~\ref{thm:dimension} is
\(\max\{0,d/\rho-(d-1)\}\), also for subexponential families
nonempty at infinitely many levels.

For two bases with \(\rho<2-\lambda\) and \(D_f>0\), the same
argument with \(a=(Nv,Mv)\) gives dimension \(D_f/\rho\) for
any subexponential family nonempty at infinitely many levels.
We next treat \(\rho>2-\lambda\). At a fixed level and for fixed
words in both bases, let \(c,c_0\) be two solutions with initial
integer vectors \(p_1,p_2\) and \(p_{1,0},p_{2,0}\), respectively.
The errors \(cv-p_1/N\) and \(cv-p_2/M\) lie in the prescribed
scaled boxes. Subtracting the errors for \(c_0\) from those for
\(c\) cancels the box offsets, so
\[
 \frac{M}{G}(p_1-p_{1,0})-\frac{N}{G}(p_2-p_{2,0})
   = O(Q^{2-\lambda-\rho}).
\]
This is an integer vector and hence vanishes at all sufficiently
large levels. The differences of the rational approximants both
equal \(p/G\) for some \(p\in\mathbb Z^d\). With \(c_0\) fixed,
we obtain
\[
  G(c-c_0)v-p=O(GQ^{-\rho}).
\]
If \(D_g>0\), apply the upper count in
Lemma~\ref{lem:met-segment} with \(H=G\), \(\beta=\lambda>0\),
and \(D=D_g\). The parameter interval for \(c-c_0\) has bounded
length, and the estimate is uniform over the choice of words.
This gives an upper bound \(D_g/\rho\) for subexponential
families. Conversely, if \(Gcv\) lies above an integer vector by
coordinatewise amounts in a box of the form
\((a_0GQ^{-\rho},a_1GQ^{-\rho})^d\) with
\(0<a_0<a_1<1/2\), both radix windows are zero. Indeed, the integer
vector on the \(G\)-grid gives integer vectors on the \(N\)- and
\(M\)-grids, and in each coordinate the error in \(cv\) is below
both \(B_1^{-\lfloor\rho u\rfloor}\) and
\(B_2^{-\lfloor\rho u_2\rfloor}\). The lower bound of
Lemma~\ref{lem:met-segment} gives dimension at least \(D_g/\rho\)
for the joint zero windows.

Since \(D_f-D_g=d(2-\lambda-\rho)\), the leading positive term
in each of these two ranges is the larger of the two. If the
terms agree at a positive value, use shorter windows for the
upper bound and longer zero windows for the lower bound, then
pass to the common exponent by continuity. At and above
\(\rho_*\), truncate the allowed words to exponents approaching
\(\rho_*\) from below. Both steps preserve the subexponential
bound on the number of words in an upper cover. This proves the
joint zero-window dimension
\[
  \frac{\max\{0,D_f,D_g\}}{\rho}.
\]
It gives the same upper bound for all subexponential families,
and equality for families which include the joint zero word
at infinitely many levels. These lower bounds hold along any
infinite subsequence of levels. In particular, one or two
windows with \(o(u)\) nonzero vector positions allowed have the
same dimension as the corresponding zero windows. To count
these patterns, if \(L=O(u)\) is the total number of positions
and \(k=o(u)\) the defect allowance, then for a fixed bound
\(A\) on the number of symbols at a position the count is at
most \(\sum_{j=0}^k\binom{L}{j}A^j=\exp(o(u))\). This follows
from \(\binom{L}{j}\leq(eL/j)^j\) for \(j>0\) and
\((k/u)\log(u/k)\to0\).

Arbitrary prescribed pairs of words can be incompatible when
\(\rho>2-\lambda\). In one coordinate take the first word zero.
Choose the starting point of the second scaled cylinder within
\(O(Q^{-\rho})\) of \(G/(2NM)\), which is possible because its
starting points have spacing \(B_2^{-\lfloor\rho u_2\rfloor}\)
and \(G/(2NM)\leq1/(2M)\). The difference of the two rational
initial coordinates lies in \((G/(NM))\mathbb Z\), whereas the
required difference of their errors is a half-step away from
this grid up to \(O(Q^{-\rho})=o(G/(NM))\). There is then no
solution for large \(u\). At \(\rho=2-\lambda\) the zero-word
lower bound follows by monotonicity as above.

The strict cutoff in one base for boundedly many nonzero vector
digits follows from Lemma~\ref{lem:transitions}, by taking
\(W_m=\lfloor cB^mv\rfloor\) coordinatewise. A subexponential
family without a restriction on individual words may have
dimension zero and still contain a fixed \(c\) at every level,
by prescribing just the observed word of that \(c\).

For the strict cutoff for joint zero windows, suppose a fixed
\(c>0\) has infinitely many levels at which the two initial
vectors satisfy
\begin{equation}\label{eq:met-joint-primal}
  p_1/N=cv+O(Q^{-\rho}),\qquad p_2/M=cv+O(Q^{-\rho}).
\end{equation}
When \(\rho>2-\lambda\), their difference is in
\((G/(NM))\mathbb Z^d\), so the two vectors agree for all
sufficiently large such levels. The following descent also
handles \(\rho\leq2-\lambda\).

Suppose \(\ell\) independent integer linear projections of the
two vectors already agree along infinitely many levels. Initially
take \(\ell=0\). Adjoin integer rows to obtain a basis of
\(\mathbb Q^d\), and call the transformed direction \(v'\).
Its coordinates are still linearly independent over \(\mathbb Q\).
The shared coordinates have denominators dividing \(G\), since
each has denominators dividing both \(N\) and \(M\). Keep a single
copy of each of these \(\ell\) coordinates and keep both copies
of each remaining coordinate. The resulting \(m=2d-\ell\)
rationals have designated denominators \(H_i=G,N,M\) in the
indicated blocks. Let \(J=\operatorname{lcm}(N,M)\) and multiply
the rational vector by \(J\) to obtain \(Y\in\mathbb Z^m\).
Its limiting direction after division by \(J\) is \(c\widetilde v\),
where \(\widetilde v\) lists the entries of \(v'\) with those
multiplicities. All entries of \(v'\) are nonzero.

At infinity use \(Y_1\) and the \(m-1\) independent differences
\(Y_i-(\widetilde v_i/\widetilde v_1)Y_1\). The real product is
\(O(J^mQ^{-(m-1)\rho})\) by \eqref{eq:met-joint-primal}. Since
\(J/H_i\) divides \(Y_i\) and is supported on the primes dividing
\(B_1B_2\), the finite-place coordinate products give the bound
\[
 O\!\left(Q^{\,2(d-\ell)+\lambda\ell-(2d-\ell-1)\rho}\right)
 = O\!\left(Q^{(1-\ell/d)D_f+(\ell/d)D_g}\right)
\]
for the full product. If \(\rho>\rho_*\), both \(D_f,D_g<0\).
As \(\|Y\|_\infty\asymp J\) grows at most polynomially in \(Q\),
the Subspace Theorem supplies a fixed nontrivial rational relation
along an infinite subsequence. Dividing by \(J\) and taking the
limit, rational independence of \(v'\) forces each shared
coefficient to vanish, and forces the two coefficients on each
duplicated coordinate to be opposites. Thus the relation gives
a further common projection, independent of the previous ones.
At \(\ell=d\) the limiting rational independence precludes any
nontrivial relation. This proves the strict cutoff, separately
from the dimension-zero conclusion at equality.

\subsection{Occupancy at the one-base endpoint}

\begin{proposition}\label{prop:met-critical-field}
Suppose \(v_1,\ldots,v_d\) form a \(\mathbb Q\)-basis of a
degree-\(d\) number field \(K\subset\mathbb R\). For every integer
\(B\geq2\), the simultaneous zero-window set at
\(\rho=d/(d-1)\) contains a dense \(G_\delta\) subset of some
half-line \((C,\infty)\), and has Hausdorff dimension zero.
\end{proposition}

\begin{proof}
The dimension follows from Theorem~\ref{thm:dimension}. Choose
one representative of each complex conjugate pair of embeddings.
Let \(V\) be the hyperplane in the resulting logarithmic
coordinates with weighted sum zero, using weight \(1\) at real
and \(2\) at complex places. Use the distinguished real embedding
as the first coordinate. By Dirichlet's unit theorem
\cite[Chapter~I, \S7]{met-neukirch}, the logarithms of absolute
values of squared units form a full lattice \(\Lambda\) in \(V\).
Put
\[
 w=\left(\log B,-\frac{\log B}{d-1},\ldots,
                  -\frac{\log B}{d-1}\right)\in V.
\]

Write \(P=\overline{\Lambda+\mathbb Zw}\) and let \(U\) be the
largest vector subspace contained in \(P\), obtained by adding all
such subspaces. The closed subgroup \(P\cap U^\perp\) is discrete.
Otherwise a sequence of its nonzero vectors tending to zero has,
after normalization and passage to a subsequence, a limiting unit
direction. Integer multiples of those vectors approximate every
real multiple of that direction, adding a line to \(U\). Since
\(U\subset P\), projection to \(U^\perp\) sends \(P\) onto
\(P\cap U^\perp\). The projection of \(\Lambda\) spans \(U^\perp\),
so a bounded fundamental parallelepiped of a basis chosen from
this projection contains only finitely many points of
\(P\cap U^\perp\). Hence \(P/(U+\Lambda)\) is finite. If the
first coordinate of \(U\) were zero, the relation \(mw\in
U+\Lambda\) for some positive integer \(m\) would give a squared
unit with distinguished value \(B^m\). Injectivity of the real
embedding would make it the rational integer \(B^m\), which
cannot be a unit. Choose \(h\in U\) with first coordinate \(1\).

In the compact torus \(V/\Lambda\), the closure of each forward
tail of \(u w\) equals the closure of the cyclic group generated by
\(w\). Indeed a convergent subsequence with increasing gaps gives,
by taking differences, positive integers \(n_j\to\infty\) with
\(n_j w\to0\) in the torus. Then \((n_j+m)w\to mw\) for every
integer \(m\). Consequently, for each \(t\in[0,\log B]\) there
are \(u_j\to\infty\) and squared units \(\epsilon_j\) such that
their log vectors minus \(u_jw\) tend to \(th\).
The distinguished values satisfy
\(\epsilon_j/B^{u_j}\to e^t\), and at every other embedding the
modulus is \(O(B^{-u_j/(d-1)})\), with an eventual common constant
for \(0\leq t\leq\log B\).

Choose an integer \(D_0>0\) such that all \(D_0v_i\) are
algebraic integers. The trace
\(p_{i,j}=\operatorname{Tr}_{K/\mathbb Q}(D_0\epsilon_jv_i)\)
is a rational integer and
\(D_0\epsilon_jv_i-p_{i,j}=O(B^{-u_j/(d-1)})\) by the bounds on
the other conjugates. Fix an integer \(k\) sufficiently large
and put \(m_j=u_j-k\), \(c_j=D_0\epsilon_j/B^{m_j}\).
The possible limits of \(c_j\) fill
\([D_0B^k,D_0B^{k+1}]\). At the level \(m_j\), the trace errors
in \(B^{m_j}c_jv_i\) are at most \(O(B^{-k/(d-1)})\) times the
critical box width \(B^{m_j-\lfloor\rho m_j\rfloor}\).
Let \(v_{\max}=\max_i v_i\) and \(v_{\min}=\min_i v_i>0\).
Increasing \(c_j\) by
\(B^{-\lfloor\rho m_j\rfloor}/(2v_{\max})\) adds between
\(v_{\min}/(2v_{\max})\) and \(1/2\) times that width in each
coordinate. For large \(k\) this places every coordinate
strictly inside the zero box. The shifted centers retain the
same limits. For each lower cutoff on \(m_j\), the union of all
such strict zero-window events is open and dense in the tail
covered by the adjacent intervals
\([D_0B^k,D_0B^{k+1}]\). Intersecting these unions gives the
claimed dense \(G_\delta\).
\end{proof}

The restriction to a tail can matter. For \(v=(1,\sqrt2)\) and
\(\rho=2\) let \(q=\lfloor B^u c\rfloor\) and
\(p=\lfloor B^u c\sqrt2\rfloor\). A critical zero window implies
\(|q\sqrt2-p|\leq\sqrt2/B^u\). At large \(u\), we have
\(p+q\sqrt2<4q\). The nonzero integer \(p^2-2q^2\) then gives
\(|q\sqrt2-p|\geq1/(4q)\), which excludes all sufficiently late
critical windows if \(0<c<1/(4\sqrt2)\).

\subsection{Empty windows for the binomial prefix test}

Let \(B_n=\binom{2n}{n}\), and let \(H_{g,\rho}(X)\) count the
integers \(X\leq n\leq X^\rho\) for which \(\lfloor gB_n\rfloor\)
begins with the full binary expansion of \(n\). This is the
multiplier version of the binary test in
Theorem~\ref{thm:headline}(i). Put \(c=(g/\sqrt\pi)^{2/3}\) and
\(\tau=\sqrt[3]{4}\). The change between \(g\) and \(c\) is a
homeomorphism and is bi-Lipschitz on compact positive intervals.

A hit means that for some integer \(k\geq0\),
\(2^kn\leq gB_n<2^k(n+1)\), since both thresholds are integers.
Put \(f=2^{-k}gB_n-n\), \(j=2n-k\), and \(x=c\tau^j\).
The two-term central-binomial expansion \cite{app-brent} gives
\[
 n+f=\frac{x^{3/2}}{\sqrt n}
       \left(1-\frac{1}{8n}+O(n^{-2})\right).
\]
Inverting with \(0\leq f<1\), we obtain
\begin{equation}\label{eq:met-bin-inverse}
 x=n+\frac{1}{12}+\frac23 f+O(n^{-1}),\qquad
 \|x\|_{\mathbb R/\mathbb Z}\geq \frac1{12}-O(n^{-1}),
\end{equation}
where the norm is distance to the nearest integer. This is the
inverse form of the aligned prefix test used in
Theorem~\ref{thm:transfer}. Conversely, for \(j=3m+r\) with
\(0\leq r\leq2\), test \(n=\lfloor x\rfloor\) and
\(k=2n-j\), where \(x=c4^m\tau^r\). On any fixed compact interval
of positive \(c\), we have \(k\geq0\) for all sufficiently large
\(m\). Expanding the same formula with \(x=n+\{x\}\) gives
\begin{equation}\label{eq:met-bin-direct}
  2^{-k}gB_n-n=\frac32\{x\}-\frac18+O(4^{-m}).
\end{equation}

\paragraph{Upper bound.}
Fix \(0<A<B\) so that \(c\in[A,B]\), and consider
\(1<\rho<3/2\). For a window starting at \(X\), take
\[
 u=\left\lceil\log_4((X+1)/A)\right\rceil,\qquad
 w=\lfloor\rho u\rfloor-L.
\]
Choose a fixed integer \(L\) sufficiently large in terms of the
compact interval and \(\rho\). The candidates
\(n=\lfloor c4^m\tau^r\rfloor\), \(u\leq m\leq w\),
\(0\leq r\leq2\), then belong to \([X,X^\rho]\) for all large
\(X\). Indeed \(c4^u\geq X+1\) and
\(c4^w\tau^2\leq B\tau^2\,4^{-L}(4(X+1)/A)^\rho\).
They are distinct for large \(X\), since the unrounded candidates
in increasing \(j\) have ratio \(\tau>1\).

For a fixed branch \(r\), let \(W_{m,r}\) be the nearest integer to
\(c4^m\tau^r\) and \(z_m=c4^m\tau^r-W_{m,r}\), with ties resolved
by choosing the larger integer. By \eqref{eq:met-bin-direct}, on a
failed test
\[
 -\frac14-K4^{-m}\ \leq\ z_m\ \leq\ \frac1{12}+K4^{-m}
\]
for a uniform \(K>0\). In a run of consecutive failures, the
integer \(d_m=z_{m+1}-4z_m\) on every internal edge therefore
belongs to \(\{0,1\}\) for large \(u\). If two successive nonzero
internal edges occur at \(m<l\), then
\[
 z_{m+1}\geq-4K4^{-m},\qquad
 z_l=4^{l-m-1}z_{m+1},\qquad z_l=(z_{l+1}-1)/4<-1/5.
\]
These bounds force \(l\geq2m-O(1)\). This is impossible for
\(m\geq u\) and \(l\leq w\) at large \(u\), and the argument
works for any fixed \(\rho<2\). If there are \(h_r\) successes
in the branch, there are at most \(h_r+1\) failure runs and
\(2h_r\) edges touching successes. Thus its integer transitions
\(W_{m+1,r}-4W_{m,r}=-d_m\) are nonzero at most \(3h_r+1\) times
in the window.

If \(H_{g,\rho}(X)\leq H\) for some fixed integer \(H\), there are
at most \(3H+3\) nonzero transitions in the vector \(W_m\) formed
by the three \(W_{m,r}\). All transition entries are uniformly
bounded by the rounding rule. For each starting level \(u\)
there are \(O_H(u^{3H+3})\) lists of nonzero positions and
values. On a fixed list, telescoping with \(v=(1,\tau,\tau^2)\)
gives
\[
 c4^u v=W_u+
   \sum_{m=u+1}^{w}4^{u-m}(W_m-4W_{m-1})+O(4^{u-w}).
\]
The sum is a fixed translation. The coordinates of \(v\) are
linearly independent since \(T^3-4\) is an irreducible cubic.
The upper count of Lemma~\ref{lem:met-segment} with \(Q=4^u\)
and \(\beta=1\) covers each list by \(O(Q^{3-2\rho})\) intervals
of length \(O(Q^{-\rho})\). The constant margin \(L\) changes
only the implicit constants. Consequently the \(s\)-power sum
over all lists and levels converges for \(s>3/\rho-2\).
Arbitrarily large real \(X\) yield infinitely many integer
starts \(u\). Taking the countable union over \(H\) gives the
upper bound for the set where the counts fail to tend to
infinity. This includes the empty-window set.

\paragraph{Lower bound.}
Fix \(\rho<\sigma<3/2\), and prescribe zero base-four digits in
all three coordinates of \(cv\) from position \(u+1\) through
\(\lfloor\sigma u\rfloor+K_0\), where \(K_0\) is fixed and
\(4^{-K_0}<1/24\). By the one-base proof, the limsup of these
conditions has local dimension \(3/\sigma-2\). At every such
level the fractional parts of all \(c4^m\tau^r\) with
\(u\leq m\leq\lfloor\sigma u\rfloor\) are smaller than
\(4^{-K_0}\). Take \(X=2c4^u\). For a hit in \([X,X^\rho]\),
\eqref{eq:met-bin-inverse} would give
\[
 c4^u < x < c4^{\lfloor\sigma u\rfloor}
\]
for all sufficiently large selected starts, because
\(\rho<\sigma\). The corresponding branch index lies in the
prescribed range, so its distance to an integer is below
\(1/24\), in contradiction with \eqref{eq:met-bin-inverse}.
These parameters have actual empty windows. Letting
\(\sigma\) decrease to \(\rho\) gives the lower bound
\(3/\rho-2\) for both sets.

\paragraph{The critical and supercritical exponents.}
For \(\rho=3/2\), the exceptional sets for bounded counts are
contained in those for every smaller exponent. The upper bounds
just proved give dimension zero. For \(\rho>3/2\) choose
\(3/2<\sigma<\min(\rho,2)\) and apply the candidate construction
with \(\sigma\) in place of \(\rho\). A bounded count along
arbitrarily late windows would give boundedly many transitions
on \([u,\lfloor\sigma u\rfloor-L]\) for
\(W_m=c4^m(1,\tau,\tau^2)+O(1)\). Restrict to
\([u,\lfloor\sigma_0u\rfloor]\) for \(3/2<\sigma_0<\sigma\).
Lemma~\ref{lem:transitions} contradicts a uniform bound on these
transitions for any fixed \(c>0\). This gives emptiness of both
exceptional sets above \(3/2\).

At the critical exponent, the set of parameters giving arbitrarily
late empty windows contains a dense \(G_\delta\) on the whole positive
half-line. To prove this, put \(a=\sqrt[3]{2}\), so \(\tau=a^2\),
and work in \(\mathbb Q(a)\). The element
\[
 U_0=(1+a+a^2)^2=5+3\tau+2\tau^2
\]
is a unit of norm \(1\), since \((a-1)(1+a+a^2)=1\).
Set \(\eta_0=1/(3\tau^2)\) and \(S_h=\eta_0U_0^h\).
For \(P\in\mathbb Z[\tau]\), the trace
\(\operatorname{Tr}(\eta_0P)\) equals the coefficient of \(\tau^2\)
after reducing \(P\) using \(\tau^3=4\). In fact
\(\eta_0=\tau/12,\ \eta_0\tau=\tau^2/12,\ \eta_0\tau^2=1/3\),
and \(\tau,\tau^2\) have trace zero. Thus
\(\operatorname{Tr}(S_h\tau^r)\) is an integer for \(r=0,1,2\).
Also \(N(\eta_0)=1/432\), so the nonreal conjugates of \(S_h\)
both have modulus \(1/\sqrt{432S_h}\). Each trace error satisfies
\[
  \big|S_h\tau^r-\operatorname{Tr}(S_h\tau^r)\big|
       \leq\frac{2\tau^r}{\sqrt{432S_h}} .
\]

For any integer \(u\geq0\) put \(c_0=S_h/4^u\) and \(X=3S_h/4\).
If \(x=c_0\tau^j\geq X\), then \(j\geq3u\) because
\(S_h/\tau<X\). Write \(j=3u+3l+r\) with \(l\geq0\) and
\(r\in\{0,1,2\}\). Scaling the trace estimate by \(4^l\) gives
\[
 \|x\|_{\mathbb R/\mathbb Z}\leq
    \frac{2x}{\sqrt{432}\,S_h^{3/2}} .
\]
For \(X\leq x\leq X^{3/2}+2\), this is at most
\(1/16+O(S_h^{-3/2})\), since
\(2(3/4)^{3/2}/\sqrt{432}=1/16\).
For all sufficiently large \(h\), this excludes any hit in
\([X,X^{3/2}]\) by \eqref{eq:met-bin-inverse}. It also excludes
hits on an open neighborhood of \(c_0\). Indeed, for \(c\)
satisfying \(|c/c_0-1|<1/(1000(X^{3/2}+2))\), a hypothetical
hit would have \(c\tau^j\in[X+1/24,X^{3/2}+1]\) by
\eqref{eq:met-bin-inverse}. Then \(c_0\tau^j\in[X,X^{3/2}+2]\)
and \(|c\tau^j-c_0\tau^j|\leq1/500\). Its distance bound is
still strictly below the necessary one for large \(h\).

Taking field norms in a putative relation \(U_0^q=4^p\) with
\(p,q>0\) gives \(1=4^{3p}\), so \(\log_4 U_0\) is irrational.
The values \(\log_4\eta_0+h\log_4 U_0\) are dense modulo \(1\)
on every forward tail. Indeed, the cyclic group generated by
an irrational rotation has nonzero residues arbitrarily close
to zero by pigeonholing, and multiples of such a residue
approximate every point of the circle. The forward tail closure
equals the group closure by the same difference argument as
in Proposition~\ref{prop:met-critical-field}. Choosing \(u\) to
account for the integer part gives centers \(S_h/4^u\) in every
positive open interval for arbitrarily large \(h\).
For each \(X_0>0\), the union of the neighborhoods just proved
empty for some \(X>X_0\) is open and dense. The intersection
over positive integer cutoffs \(X_0\) is the claimed dense
\(G_\delta\). The bi-Lipschitz change on compact intervals
transfers the dimension conclusions to \(g\), and its
homeomorphism transfers the dense \(G_\delta\).

\subsection{Relation to earlier work}

The ambient restricted-denominator dimension problem goes back to
Borosh and Fraenkel \cite{met-borosh-fraenkel}, with a weighted
extension by Rynne \cite{met-rynne}. For affine subspaces,
Huang and Liu \cite[Theorem~7]{met-huang-liu} give a fixed-denominator
count, and Huang \cite[Theorems~3.1(b) and~3.2(b)]{met-huang}
gives a shifted upper count. For an algebraic line the latter
already yields the one-base upper count in the positive-dimension
range. Baier and Ghosh \cite{met-baier-ghosh} study a different
restriction on a line through the origin, imposing primality on
the denominator and one numerator. Dynamical dimension results
for independently varying coordinates include Wu's simultaneous
\(\beta\)-system theorem \cite{met-wu}. Bugeaud and Liao
\cite{met-bugeaud-liao} treat uniform approximation by radix powers,
and Bugeaud and Kim \cite{met-bugeaud-kim} treat block complexity
of two expansions of one real number.

Li, Liao, Velani, and Zorin discuss two shrinking targets on a
bi-Lipschitz planar graph \cite[Eq.~(75)]{met-li-liao}. On the
scalar diagonal they report dimension \((1-\theta)/(1+\theta)\)
for the maximum of the errors at \(2^n\) and \(3^n\) bounded by
\(3^{-n\theta}\) infinitely often, for \(0\leq\theta\leq1\).
Their proof is conditional on \(abc\) when
\(\theta>1-\log 2/\log 3\) \cite[Remark~13]{met-li-liao}.
Li, Li, and Wu \cite[Theorem~1.3 and \S4]{met-li-li-wu} obtain a
general scalar maximum-norm convergence bound with a denominator-size
term and a gcd term. For equal-index powers \(a^n,b^n\), with \(b>a\),
and errors of size \(b^{-n\theta}\) for \(\theta>1\), their exact
dimension formula is \(\log_b\gcd(a,b)/(1+\theta)\).
Zhang and L\"u \cite{met-zhang-lu} treat the product of errors
for equal-index powers of \(2\) and \(3\). Tapia
\cite[Theorem~1]{met-tapia} gives a gcd-sensitive convergence
criterion for inhomogeneous products with integer denominators. The additional feature of Theorem~\ref{thm:dimension}
is the exact competition between the two terms on a line with
\(d\geq2\) rationally independent algebraic coordinates in its
direction and comparable physical starting scales, with one shared
scalar parameter.

\section{Prefixes and suffixes of polynomial values}\label{sec:polynomials}

\subsection{Unshifted polynomial prefixes}\label{subsec:polynomial-prefixes}

Let \(P\in\mathbb Q[X]\) be eventually integer-valued, with
\(P(X)=aX^{q+1}+A X^q+O(X^{q-1})\), \(a>0\), and \(q\ge2\).
For a trial suffix length \(k\), write
\[
 Q=b^k,\qquad x=(Q/a)^{1/q},\qquad t=A/(qa).
\]
If \(P(n)=Q(n+u)\) with \(0\le u<1\), division by \(n+u\)
gives \(x^q=n^q(1+O(1/n))\), hence \(n=x+O(1)\).
Uniformly in that range,
\begin{equation}\label{eq:pol-prefix-residual}
 P(n)/Q-n=q(n-x+t)+O(x^{-1}).
\end{equation}
In the branch \(k=qm+r\), \(0\le r<q\), put
\(\beta_r=(b^r/a)^{1/q}\), so \(x=\beta_r b^m\). If \(\beta_r\)
is irrational then, for all sufficiently large \(m\), a prefix
match occurs precisely when
\begin{equation}\label{eq:pol-irrational-arc}
 x\bmod 1\in (t-1/q,t)\bmod 1,
 \qquad n=\lceil x-t\rceil\ \text{on a match}.
\end{equation}
Indeed, being within \(O(x^{-1})\) of either rational endpoint
infinitely often would approximate \(\beta_r\) to
\(O(b^{-2m})\) by rationals with reduced denominators dividing
\(M b^m\), for a fixed common denominator \(M\).
These denominators are unbounded along any such sequence and use
only a fixed set of primes. Ridout's theorem~\cite{pol-ridout}
excludes such approximations, since for any fixed \(0<\eta<1\)
the error is eventually smaller than the reduced denominator
to the power \(-1-\eta\). Equation~\eqref{eq:pol-prefix-residual}
then gives the exact test. If \(\beta_r\) is rational, the
fractional parts of \(x\) follow an eventual finite cycle. For
each cycle value there is at most one integer \(n\) with
\(0\le n-x+t\le1/q\), and the gap to all other candidates is
bounded below. For the possible candidate, substitute its fixed
offset \(n-x\) into \(P(n)-a x^q n\) and \(P(n)-a x^q(n+1)\).
The eventual signs decide the weak lower and strict upper
inequalities, including the endpoints.

\begin{proposition}\label{prop:pol-power-prefix}
For \(P(n)=n^d+c\), \(c\in\mathbb Z\), \(d\ge2\), the following
classifies infinitude of prefix matches:
\begin{enumerate}
\item If \(c\ge0\), every radix has infinitely many matches.
\item If \(d=2\) and \(c<0\), every radix has infinitely many
matches for \(c\le-2\). For \(c=-1\) there are no matches.
\item If \(d\ge3\) and \(c<0\), write \(q=d-1\).
Every perfect \(q\)-th power radix gives finitely many matches,
and none for \(c=-1\). In any other radix there are infinitely
many matches if and only if at least one irrational branch
\(b^{r/q}b^m\), \(0\le r<q\), visits \((1-1/q,1)\) modulo one
infinitely often. This criterion is independent of the
negative value of \(c\).
\end{enumerate}
\end{proposition}
\begin{proof}
For \(c\ge0\), take \(n=b^m\) and \(Q=b^{(d-1)m}>c\).
For \(d=2\), \(c<0\), any match requires \(n>Q\), and the
residual \(n(n-Q)+c\) at \(n=Q+1\) equals \(Q+1+c\).
It lies in \([0,Q)\) for large \(Q\) if \(c\le-2\).
For \(c=-1\) it is already \(Q\) there and increases for
larger \(n\).

For \(q\ge2\), every rational \(x=b^{k/q}\) is integral.
At a match the lower bound forces \(n\ge x+1\), so
\[
 (n+1)Q\le(n^2-1)(n-1)^{q-1}
 \le n^{q+1}-n^{q-1}\le n^{q+1}+c
\]
when \(n^{q-1}\ge |c|\), contradicting the strict upper bound.
Thus a rational scale gives no large matches, and none at
all for \(c=-1\). For irrational scales apply
\eqref{eq:pol-irrational-arc} with \(t=0\).
There are finitely many candidates at each fixed \(k\) since
\(P(n)/n\to\infty\), and the prefix intervals for different
\(k\) at a fixed \(n\) are disjoint.
\end{proof}

For \(d=3\) and \(c<0\) the test simplifies for \(n>|c|\) to
\(n^2-n+1\le Q<n^2\): the lower prefix bound gives \(Q<n^2\),
and \(P(n)/(n+1)=n^2-n+1+(c-1)/(n+1)\) lies strictly between
the two consecutive integers \(n^2-n\) and \(n^2-n+1\).
Equivalently \(n-\tfrac12<\sqrt Q<n\). For \(c=-1\) this covers
every possible matching index. In radix \(2\) the odd lengths
give scales \(2^m\sqrt2\), and the digit \(1\) recurs in every
irrational binary expansion. In radix \(3\), infinitude is
equivalent to recurrence of the ternary digit \(2\) in
\(\sqrt3\). A tail beginning with \(2\) exceeds
\(1/2=0.\overline{1}_3\), whereas a tail consisting only of \(0\)
and \(1\) cannot exceed it. This entry remains open. An eventual
tail over \(0,1\) would yield a quadratic irrational in the
middle-third Cantor set upon doubling that tail, a special
case of the issue raised by Mahler~\cite[\S2]{pol-mahler}.
The conjecture that every finite word recurs in each irrational
algebraic expansion would give infinitude in precisely the
non-perfect-\(q\)-th-power radices in part~(iii).
For example, when \(d=5\) and \(b=4\), the odd lengths give
\(2^m\sqrt2\) and the arc is \((3/4,1)\). This entry asks
whether \(11\) recurs in the binary expansion of \(\sqrt2\).

\begin{proposition}\label{prop:pol-pascal-prefix}
For \(d=3,5\) put \(q=d-1\), \(H=d!\), and \(x=(H b^k)^{1/q}\).
The radices admitting an integral scale are exactly
\(b\in 6(\mathbb Q_{>0})^2\) for \(d=3\), and
\(b\in 120(\mathbb Q_{>0})^4\cup120^3(\mathbb Q_{>0})^4\)
for \(d=5\), in each case restricted to integer radices.
These scales never give a match. On the remaining lengths the
eventual arcs are \((0,1/2)\) for \(d=3\) and \((1/4,1/2)\)
for \(d=5\), with \(n=\lceil x+d/2\rceil\). The cubic arc test
holds at every length. In particular, \(\binom n3\) has
infinitely many matches in radices \(2,4\), and \(\binom n5\)
has infinitely many in radices \(2,4,16\).
\end{proposition}
\begin{proof}
The coefficient ratio gives \(t=-d/2\).
An integral scale occurs exactly when \(k v_p(b)+v_p(H)\)
is divisible by \(q\) at every prime \(p\), where \(v_p\) is
the usual prime valuation. A prime with \(v_p(H)=1\) forces
\(k\) to be a unit modulo \(q\). The congruences are thus
equivalent to \(b/H^u\) being a rational \(q\)-th power,
with \(u\equiv-k^{-1}\pmod q\). Thus for each listed radix
the integral scales are precisely the lengths in the class
\(k\equiv-u^{-1}\pmod q\). The possible units \(u\) are \(1\)
for \(q=2\) and \(1,3\) for \(q=4\).

For \(d=3\), the exact inequalities at \(n\ge3\) read
\[
 (n-2)^2+2-\frac{6}{n+1}<6Q\le(n-1)(n-2).
\]
The integer \((n-2)^2+1\) is never divisible by \(3\), so
these are equivalent to \(n-2<x<n-\tfrac32\). This excludes
integral \(x\) and gives the exact arc. For \(d=5\), the
arithmetic and geometric mean inequality on \(n-1,\ldots,n-4\)
forces \(n>x+5/2\) at a match. If \(x\) is integral, it is
divisible by \(30\) and \(n\ge x+3\). At \(n=x+3\),
\[
 \frac{120\binom n5}{n+1}
 =\frac{x+3}{x+4}\,x(x+1)(x-1)(x+2)>x^4.
\]
The expression on the left increases for \(n\ge5\), so
the upper prefix inequality fails. The eventual arc at the
irrational scales follows from~\eqref{eq:pol-irrational-arc}.

The listed radices give a branch \(x=2^m H^{1/q}\), by taking
\(k=qm/\log_2 b\). The fractional part of \(H^{1/q}\)
is irrational. The cubic arc reads a binary \(0\), which
recurs, and the quintic arc reads \(01\), whose eventual
absence would force an eventually constant binary tail.
\end{proof}

These positive cases leave other radices subject to their
irrational-branch tests. For example, for \(\binom n3\) in
radix \(9\), infinitude is equivalent to recurrence of the
ternary digit \(0\) in \(\sqrt6\), using
\(1/2=0.\overline{1}_3\). In radix \(3\) the two branches give
the same question for at least one of \(\sqrt6,\sqrt2\).

\begin{proposition}\label{prop:pol-necklace-complement}
Fix \(h\ge1\), and let \(L_{b^h}(n)\) count primitive necklaces
of length \(n\) on \(b^h\) letters. For nonsquare \(b\), the
prefix matches of \(n^3-1\) and of \(L_{b^h}(n)\) eventually
partition the nearest integers to \(b^m\sqrt b\), \(m\ge0\).
Their combined count up to \(X\) is \(\log_b X+O_{b,h}(1)\).
For square \(b\) both match sets are finite.
\end{proposition}
\begin{proof}
Each primitive necklace has \(n\) rotations. An imprimitive
word repeats a block of length at most \(n/2\), and constant
words are imprimitive for \(n\ge2\). Thus
\(nL_{b^h}(n)=b^{hn}(1-\epsilon_n)\), where
\(0<\epsilon_n\le n b^{-hn/2}\).
The prefix bounds are \(n^2\le T(1-\epsilon_n)<n(n+1)\)
for \(T=b^{hn-k}\). Necessarily \(T>n^2\), so \(T\) is integral.
For large \(n\), \(\epsilon_n(n(n+1)+1)<1\). Consequently
the two bounds are equivalent to \(n^2<T\le n(n+1)\).
This interval contains no square, so when \(b\) is nonsquare
the only possible exponents are \(hn-k=2m+1\).
Integrality of \(T\) makes the test equivalent to
\(n<b^m\sqrt b<n+\tfrac12\), and the resulting \(k\) is
nonnegative for large \(m\). The cube test gives the other
half, \(n-\tfrac12<b^m\sqrt b<n\). The scales are irrational,
so there are no ties. They grow by a factor of \(b\), which
gives the count and eventual disjointness of the indices.
In a square radix all scales are integral, so neither test
gives a large match.
\end{proof}

\subsection{Deciding simultaneous prefix and suffix matches}
\label{subsec:polynomial-joint}

Let \(\ell_b(n)\) be the number of digits of \(n\) in radix \(b\).
A joint match for \(P\) satisfies the prefix inequalities and
\(P(n)\equiv n\pmod{b^{\ell_b(n)}}\).

\begin{theorem}\label{thm:pol-joint}
Let \(P\in\mathbb Q[X]\) be eventually integer-valued, of
degree \(D\ge2\) with leading coefficient \(a>0\). Set \(q=D-1\),
\(F(X)=P(X)-X\), and \(\beta_r=(b^r/a)^{1/q}\) for
\(0\le r<q\). Apart from finitely many exceptions, a joint
match at \(k=qm+r\) has
\begin{equation}\label{eq:pol-joint-families}
 n=\beta_r b^m+\gamma,\qquad
 \beta_r,\gamma\in\mathbb Q,\quad F(\gamma)=0.
\end{equation}
A finite calculation decides whether there are infinitely
many joint matches. For each family in
\eqref{eq:pol-joint-families}, the successful sufficiently large
values of \(m\) form a computable union of residue classes.
\end{theorem}
\begin{proof}
The prefix gives \(b^k=a n^q+O(n^{q-1})\), so
\(n=\beta_r b^m+O(1)\) and \(\ell_b(n)=m+O(1)\).
At one fixed prime \(p\mid b\), the suffix gives
\(|F(n)|_p\le p^{-\ell_b(n)v_p(b)}\). Factoring \(F\) over
\(\overline{\mathbb Q_p}\) and taking the smallest of the
\(D\) root distances gives a root \(\alpha\) with
\[
 |n-\alpha|_p\le |a|_p^{-1/D}p^{-\ell_b(n)v_p(b)/D}.
\]
For each fixed \(r,\alpha\), use the following independent
rows of linear forms on the primitive vector \(z=(b^m,n,1)\):
\[
 \begin{array}{c|ccc}
  \infty & X & Y-\beta_r X & Z\\
  p & X & Y-\alpha Z & Z\\
  p'\mid b,\ p'\ne p & X & Y & Z .
 \end{array}
\]
These forms and the cancellation of the \(X\)-factors are
those used by Adamczewski and Bugeaud in their proof of
Theorem~2~\cite[\S5]{pol-sym}. Here the polynomial congruence
supplies the local coefficient \(\alpha\). The \(X\)-factors
have product one, the real difference is bounded, and all
remaining finite factors except \(Y-\alpha Z\) have norm at
most one. The full product is \(O(p^{-\ell_b(n)v_p(b)/D})\),
a negative power of \(\|z\|_\infty\asymp b^m\).
The \(S\)-adic Subspace Theorem~\cite{schlickewei}
in the integer-product form used in Section~\ref{sec:transfer}
therefore places the vectors in finitely many proper rational
subspaces. The choices of \(r\) and \(\alpha\) are finite.

On a subspace containing an unbounded sequence of \(m\),
take a fixed rational relation \(U b^m+V n+W=0\). We have
\(V\ne0\). Dividing by \(b^m\) and taking the real limit
gives \(\beta_r=-U/V\), and substitution gives
\(\gamma=-W/V\). At \(p\), \(b^m\to0\) and \(F(n)\to0\),
so \(F(\gamma)=0\) by continuity. Bounded \(m\) gives
only finitely many prefix matches. This proves the restriction.

To perform the decision, enumerate the rational \(\beta_r\)
and rational roots \(\gamma\) of \(F\). For each pair, use
the following three tests.
\begin{enumerate}
\item For a common denominator \(M\) of \(\beta_r,\gamma\),
integrality of \(n\) is the condition
\(M\beta_r b^m+M\gamma\equiv0\pmod M\).
Powers of \(b\) modulo \(M\) follow an eventual finite cycle.
\item Put \(N(T)=\beta_r T+\gamma\). The two prefix tests
are the eventual signs of
\(P(N(T))-b^r T^qN(T)\) and \(P(N(T))-b^r T^q(N(T)+1)\).
As \(T=b^m\to\infty\), these must be nonnegative and
strictly negative, respectively. The leading nonzero
coefficient determines the sign, with sign zero for the
zero polynomial.
\item Eventually \(\ell_b(n)=m+J\). If
\(b^{J-1}<\beta_r<b^J\), this determines \(J\). At
\(\beta_r=b^j\) use \(J=j\) for \(\gamma<0\) and \(J=j+1\)
otherwise. Let \(h\) be the multiplicity of \(\gamma\)
and \(C_h\) the coefficient of \(u^h\) in \(F(\gamma+u)\).
For every \(p\mid b\) that term has uniquely smallest
valuation for large \(m\), so
\[
 v_p(F(n))=hm\,v_p(b)+v_p(C_h\beta_r^h).
\]
The suffix test eventually always passes for \(h\ge2\).
For \(h=1\) it passes exactly when
\(v_p(F'(\gamma)\beta_r)\ge Jv_p(b)\) for all \(p\mid b\).
\end{enumerate}
Thus infinitude is equivalent to passing the sign and
valuation tests on a family whose integrality cycle
contains a success.
\end{proof}

The decision determines the eventual residue classes on
the rational families. The Subspace Theorem argument gives
no effective cutoff for the exceptional matches outside them.

\begin{corollary}\label{cor:pol-joint-powers}
For every fixed \(b\ge2\), \(d\ge2\) and \(c\in\mathbb Z\),
the polynomial \(n^d+c\) has only finitely many joint matches.
For \(n^2\) the only joint match is \(n=1\).
\end{corollary}
\begin{proof}
On a family from Theorem~\ref{thm:pol-joint}, a rational
\(\beta_r=b^{r/(d-1)}\) is integral, so integrality of \(n\)
forces \(\gamma\in\mathbb Z\). Put \(x=\beta_r b^m\).
The prefix residual at \(n=x+\gamma\) tends to
\((d-1)\gamma\in[0,1]\). Thus \(\gamma=0\) except
possibly when \(d=2,\gamma=1\). In that case \(F(1)=0\)
gives \(c=0\), and \((x+1)^2>(x+2)x\) violates the
upper bound. If \(\gamma=0\) then \(F(0)=0\) gives \(c=0\).
For large \(m\), each \(p\mid b\) divides \(n=x\) and
\(v_p(n^d-n)=v_p(n)\). The suffix would force
\(b^{\ell_b(n)}\mid n\), impossible for \(n>0\) of that length.
For \(n^2\), the prefix inequalities force \(n=b^k\)
at every match. When \(k\ge1\) the same valuation argument
excludes it, and \(n=1\) is a joint match.
\end{proof}

For \(n^2+c\) the finite set is computable by restricting
to \(b^k\le |c|+1\). Indeed, for \(Q=b^k>|c|+1\), the
prefix condition \(0\le n(n-Q)+c<Q\) forces
\(n\in\{Q,Q+1\}\). On either index \(P(n)-n\equiv c\pmod Q\),
and \(Q\mid b^{\ell_b(n)}\). The suffix forces \(c=0\), and
then the sole prefix candidate \(n=Q\) fails the suffix.

\begin{example}\label{ex:pol-joint-quadratic}
The polynomial \(P(n)=n^2-n+1\) has infinitely many joint
matches in every radix \(b\). For \(T=b^m\), \(m\ge1\),
take \(n=T+1\). Then \(P(n)/T=n+1/T\) and \(P(n)-n=T^2\)
is divisible by \(b^{m+1}=b^{\ell_b(n)}\).
\end{example}

Constants and eventually nonpositive polynomials have only
finitely many joint matches. For an eventually integer-valued
linear polynomial \(P(n)=an+c\) with \(a>0\), both coefficients
are integers. Infinitude of prefix matches requires \(a=b^k\)
and \(0\le c<a\), since the possible \(k\) are bounded and
the prefix residual is \((a-b^k)n+c\). For \(a=1\) this
gives \(P(n)=n\), with every index a joint match. For
\(a=b^k>1\), \(c=0\) fails the suffix because \(a-1\) is
coprime to \(b\). If \(c=a-1\), then \(n=b^\ell-1\) works
for every \(\ell\ge1\). For \(0<c<a-1\), choose a cyclic
rotation \(j\) of the \(k\)-digit padded word of \(c\) with
nonzero leading digit. Multiplication by \(b\) modulo \(a-1\)
rotates the word, so \(jb^\ell\equiv c\pmod{a-1}\) for
infinitely many \(\ell\). The index
\(n=(jb^\ell-c)/(a-1)\) then has \(\ell\) digits for large
\(\ell\), since \(1/b<j/(a-1)<1\), and satisfies the suffix
as well as the prefix.

\section{Digit reversal and Fibonacci representations}\label{sec:representations}

\subsection{Reversals of full length}
Let \(R_b(n)\) denote the integer obtained by reversing the base-\(b\)
digits of \(n\). We require \(b\nmid n\) when asking for a reversal
of full length. Otherwise \(n=b^t\) already gives the fixed reversed
prefix \(1\) for every power \(n^d\).

\begin{proposition}\label{prop:reversal}
For every sufficiently large \(L\) there is a nonpalindromic binary
integer \(n\) of length \(L\) such that \(n^2\) begins with the entire
word for \(R_2(n)\). In base four, \((2\cdot4^t+1)^4\) begins with
\(4^t+2=R_4(2\cdot4^t+1)\) for all sufficiently large \(t\).
\end{proposition}
\begin{proof}
Put \(N=2^L\), \(E=A2^t\), and \(n=N-1-E\).
Choose \(A=3\) for odd \(L\) and \(A=5\) for even \(L\), and set
\(L=2t+a+1\), where \(a\) is the binary length of \(A\).
The word for \(E\), padded to length \(L\), consists of \(t+1\) zeros,
the palindromic word for \(A\), and \(t\) zeros. Its reversal represents
\(2E\). Complementing all digits commutes with reversal, so
\(R_2(n)=N-1-2E\). Both ends of \(n\) are ones for \(t\ge1\), and the
reversal differs from \(n\). Now
\[
 \left\lfloor n^2/N\right\rfloor
 =N-2(E+1)+\left\lfloor(E+1)^2/N\right\rfloor.
\]
The last quotient tends to \(9/8\) or \(25/16\); hence its floor is
one in both parities. In base four, putting \(u=4^t\) gives
\((2u+1)^4/(16u^3)=u+2+3/(2u)+1/(2u^2)+1/(16u^3)\).
Its integer part is \(u+2\) for large \(t\), as required.
\end{proof}

\begin{proposition}\label{prop:palindromes}
Fix \(b,d\ge2\). Among the positive integers \(n\le X\) for which
\(n^d\) begins with \(n\) in base \(b\), the number of palindromes
is \(o(\log\log X)\).
\end{proposition}
\begin{proof}
Write \(q=d-1\), \(k=qm+r\), \(0\le r<q\), for the number of digits
removed at the right. Set \(x=b^{k/q}\). The prefix inequality implies
\(n\ge x\). If \(n\ge x+1\), then
\((n+1)x^q\le(n+1)(n-1)^q<n^{q+1}\), contrary to the strict upper
inequality. Thus \(n=\lceil x\rceil\). When \(x\) is integral and
\(k\ge q\), \(b\mid x\), as is seen from the prime valuations.
Such an \(n\) cannot be a palindrome.

On any remaining irrational branch,
\(\zeta=b^{r/q-1}\) lies strictly between \(1/b\) and \(1\).
Write its base-\(b\) digits as \(0.a_1a_2\ldots\).
Eventually \(L=m+1\) is the length of \(n\) and
\(n-1=\lfloor b^L\zeta\rfloor\). For large \(m\) the leading digit
of \(n\) is \(a_1>0\), since \(b\zeta\) is irrational and adding
\(b^{-m}\) cannot cross the next integer. If \(n\) is palindromic,
its last digit is also \(a_1\). Subtracting one changes only that
digit. Consequently \(a_2\ldots a_{L-1}\) is a palindromic prefix of
the digit word for the algebraic irrational \(\{b\zeta\}\).

Suppose two such prefix lengths \(u<v\) tend to infinity with
\(v/u\le C\) for a fixed \(C>1\). The shorter palindrome is both
a prefix and a suffix of the longer, giving the latter period \(v-u\).
If these periods stay bounded on a subsequence, one fixed period holds
for arbitrarily long prefixes, contradicting irrationality. Otherwise
take a subsequence with periods increasing to infinity. The digit word
then starts with at least \(v/(v-u)\ge C/(C-1)>1\) copies, allowing
a final partial copy, of its block of length \(v-u\). The stammering
criterion of Adamczewski and Bugeaud excludes this for algebraic
irrational numbers in an integer base
\cite[Theorem~5]{rev-complexity}; see also their symmetric-pattern
criterion \cite[Theorem~4 and \S6]{pol-sym}.

On each infinite branch successive such lengths therefore have ratio
tending to infinity. Above a branch-dependent cutoff they increase
by at least any prescribed factor \(K>1\). Since \(L\ll\log X\),
their count is at most \(O_K(1)+\log\log X/\log K\) on each branch.
There are only \(q\) branches. Letting \(K\to\infty\) proves the claim.
\end{proof}

The count permits infinitely many very sparse hits on an irrational
branch. The classification of full-length nonpalindromic reversal
prefixes for general pairs \((b,d)\) also remains open here. For
\(\binom{2n}{n}\), even in binary, we do not determine whether
infinitely many indices occur as both prefix and suffix, or whether
infinitely many palindromic indices occur as full prefixes.

\subsection{Quadratic prefixes in Zeckendorf numeration}\label{sec:zeckendorf}

Use \(F_0=0,\ F_1=1\), with the Fibonacci recurrence. The
Zeckendorf word of \(n>0\) represents \(n=\sum_{j\ge2}e_j F_j\),
where \(e_j\in\{0,1\}\) and adjacent digits cannot both be \(1\).
Read the digits from the highest occupied place down to \(F_2\).
We call \(n\) a hit of \(P\) if \(P(n)\) is a positive integer whose
word begins with the entire word of \(n\). The suffix length \(k\ge0\)
is the difference of the two word lengths. Put
\[
 \phi=\tfrac{1+\sqrt5}{2},\qquad \psi=-\phi^{-1},\qquad
 L_k=\phi^k+\psi^k\quad(k\ge0),
\]
so that \(L_0=2\) and \(L_1=1\). Chang and Miller use consecutive
zero-padded endpoints for fixed leading blocks and study a fixed
six-digit block of squares~\cite[Lemma~4.4, proof, and Example~4.11]{zec-chang-miller}.

\begin{theorem}[Quadratic Zeckendorf classification]\label{thm:zeck-quadratic}
Let \(P(n)=(An^2+Bn+C)/D\) be eventually integer-valued, with
\(A,D\) positive integers and \(B,C\) integers. There are computable
integers \(N,K\ge1\), an even modulus \(m\ge2\), and a finite computable
set \(\mathcal E\subset\mathbb Z\times\{0,\ldots,m-1\}\) such that for
\(n\ge N\) the hits and their suffix lengths are exactly
\begin{equation}\label{eq:zeck-lucas-families}
 n=\frac{D L_k+s}{A},\qquad
 k\ge K,\qquad k\equiv i\pmod m,\qquad (s,i)\in\mathcal E .
\end{equation}
In particular, infinitude is decidable. The number of hits up to \(Y\)
is \(\lvert\mathcal E\rvert\log Y/(m\log\phi)+O_P(1)\).
\end{theorem}

\begin{proof}
Write \(S_k(n)=\sum_{j\ge2}e_j F_{j+k}\), the value obtained by
appending \(k\) zeros. The legal strings on places \(F_2,\ldots,F_{j-1}\)
give each integer from \(0\) to \(F_j-1\) exactly once. This follows
by induction, splitting according to the digit at \(F_{j-1}\).
A zero gives \([0,F_{j-1})\), and a one forces the next digit to
zero and gives \([F_{j-1},F_{j-1}+F_{j-2})\). The empty string and
the single place \(F_2\) start the induction. The same bound on
lower places shows that the highest differing digit of two strings
determines their order, padding on the left with zeros as needed.
Deleting the last \(k\) places is therefore nondecreasing in the
integer order, and the first integer whose shortened string represents
\(n>0\) is \(S_k(n)\). Thus a hit with suffix length \(k\) is equivalent to
\begin{equation}\label{eq:zeck-cylinder}
 S_k(n)\le P(n)<S_k(n+1).
\end{equation}
This also applies when \(n+1\) has a longer word than \(n\).

Let \(z(n)=\sum_{j\ge2}e_j\psi^j\). Summing the odd and even powers
separately bounds \(z(n)\) strictly between \(-\phi^{-2}\) and
\(\phi^{-1}\). The identity \(F_{j+1}-\phi F_j=\psi^j\) gives
\begin{equation}\label{eq:zeck-shift}
 S_1(n)=\lfloor\phi n+\phi^{-1}\rfloor,\qquad
 S_k(n)=F_{k-1}n+F_k S_1(n)=\phi^k n+F_k z(n)
 \quad(k\ge1).
\end{equation}
The formula for \(k\ge1\) follows from the Fibonacci recurrence.
For the floor formula, see also the Fibonacci shift
in Northshield~\cite[\S6]{zec-northshield}.

At a hit put \(X=\phi^k\). Since \(F_k\le X\) for \(k\ge1\),
\eqref{eq:zeck-cylinder} gives \(\lvert P(n)-Xn\rvert\le2X\).
The bound holds for \(k=0\) as well since then \(P(n)=n\).
With \(Q=|B|+|C|\), for \(n\ge 4+\lceil 2Q/A\rceil\) we have
\[
 \frac{A n^2}{2D}\le P(n)\le\frac{2A}{D}n(n-2).
\]
Thus \(X/n\le2A/D\) and \(n\le4DX/A\) at such a hit. In particular,
\[
 |An-DL_k|\le |B|+|C|/n+2DX/n+D|\psi|^k
 \le M,\qquad M=Q+4A+D.
\]
It suffices to enumerate the integers \(|s|\le M\).

The pair \((L_k,L_{k+1})\bmod A\) is purely periodic, since
\((a,b)\mapsto(b,a+b)\) is invertible modulo \(A\). An even period
\(m\) is found by iterating from \((2,1)\bmod A\) until the pair
returns, then taking the least common multiple of the period and \(2\).
Fix \(s\) and \(i\) and take \(k\equiv i\pmod m\). Retain for testing
only integral \(n=(D L_k+s)/A\), an integrality test independent of \(k\)
on the class. Eventual integrality of \(P\) implies integrality at every
integer since its numerator is periodic modulo \(D\).

Let \(r\in\{0,\ldots,A-1\}\) be the residue of \(D L_{k+1}\), and
\(\varepsilon=(-1)^i\). For \(j=0,1\), use
\(\phi L_k-L_{k+1}=\sqrt5\,\psi^k\) to write
\[
 \phi(n+j)+\phi^{-1}
 =\frac{D L_{k+1}-r}{A}+\theta_j+
   \frac{\varepsilon D\sqrt5}{AX},\qquad
 \theta_j=\frac{r+\phi(s+Aj)}{A}+\phi^{-1}.
\]
For large \(k\), the floor of the last two terms is \(\lfloor\theta_j\rfloor\),
decreased by \(1\) exactly when \(\theta_j\) is integral and
\(\varepsilon=-1\). Denote this integer by \(h_j\). The rule is
valid once \(D\sqrt5/(AX)\) is less than the distance from
\(\theta_j\) to the nearest integer distinct from it. Consequently
\(S_1(n+j)=(D L_{k+1}-r)/A+h_j\).

Substituting in \eqref{eq:zeck-shift} makes both differences
\(P(n)-S_k(n+j)\), \(j=0,1\), Laurent polynomials in \(X\) over
\(\mathbb Q(\sqrt5)\). Indeed \(n=(D(X+\varepsilon/X)+s)/A\),
and \(F_k,F_{k-1},L_{k+1}\) are linear combinations of \(X\)
and \(X^{-1}\) by Binet's formula. Each nonzero polynomial has
the sign of its highest-degree nonzero coefficient once \(X>1\)
and \(X\) exceeds the sum of the absolute values of all lower
coefficients divided by the absolute value of that coefficient.
An identically zero difference meets the lower inequality and
fails the strict upper inequality. Exact comparisons in
\(\mathbb Q(\sqrt5)\) therefore determine whether to include
\((s,i)\) in \(\mathcal E\).

Take \(K\ge1\) large enough that all the computed floor and sign
bounds hold and \(D(\phi^K-1)>M\), ensuring positivity of each
candidate. Then take \(N\ge 4+\lceil 2Q/A\rceil\) with
\(N>4D\phi^K/A\), so every hit at \(n\ge N\) has \(k\ge K\).
These bounds are found by exact comparison and finite search.
Each retained residue class contributes
\(\log Y/(m\log\phi)+O_P(1)\) indices since
\(n=(D/A)\phi^k+O_P(1)\). There is no double counting for \(n\ge N\),
since \(n\) and \(P(n)\) determine \(k\), and \(n,k\) determine \(s\).
\end{proof}

\begin{corollary}[Shifted squares]\label{cor:zeck-square}
For fixed \(c\in\mathbb Z\), the sufficiently large hits of \(n^2+c\)
are exactly \(n=L_k\) with \(c+2(-1)^k\ge0\), with suffix length \(k\).
For every \(k\ge6\), \(n=L_k\) is a hit with suffix length \(k\)
if and only if \(0\le c+2(-1)^k<F_{k+2}\).
There are infinitely many hits exactly when \(c\ge-2\). The eventual
family uses only even \(k\) for \(-2\le c<2\) and both parities
for \(c\ge2\).
\end{corollary}

\begin{proof}
At a sufficiently large hit, \(s=n-L_k\) is bounded and
\((n^2+c)/X-n=s+O_c(X^{-1})\) with \(X=\phi^k\).
By \eqref{eq:zeck-cylinder} and \eqref{eq:zeck-shift}, this quantity
lies between \(-\phi^{-2}\) and \(1+\phi^{-1}\).
These bounds leave only \(s=0,1\) as \(k\to\infty\).

The identities \(L_k=F_{k+1}+F_{k-1}\) and
\(L_k^2=L_{2k}+2(-1)^k\) follow from Binet's formula
(see also~\cite[\S2, (11) and Lemma~2.2]{zec-stoll}).
For \(k\ge6\), adding \(1=F_2\) or \(2=F_3\) to \(L_k\) preserves
the separation of occupied places. Hence
\[
 S_k(L_k)=L_{2k},\qquad S_k(L_k+1)=L_{2k}+F_{k+2},\qquad
 S_k(L_k+2)=L_{2k}+F_{k+3}.
\]
For \(s=0\), the excess over the first endpoint is \(c+2(-1)^k\),
and the gap \(F_{k+2}\) tends to infinity. For \(s=1\), the excess
over \(L_{2k}\) is \(2L_k+c+1+2(-1)^k\), which eventually exceeds
\(F_{k+3}\) since \(2>\phi^3/\sqrt5\). The latter is the excluded
upper endpoint after subtracting \(L_{2k}\), so this branch gives
no sufficiently large hits.
\end{proof}

\section{Complementary regimes and open questions}\label{sec:complements}

Other growth laws lead to different index sets. We record the precise
scope of four comparisons. The linked companion notes give the proofs.

\paragraph{Euler products.}
Fix \(b\ge2,\ q=b^h,\ h\ge1\), and let \(E_q=\prod_{i\ge1}(1-q^{-i})\).
For the Gaussian binomial \(\genfrac{[}{]}{0pt}{}{2n}{n}_{\!q}\) and the Dyck-path
area polynomial \(\sum_P q^{\operatorname{area}(P)}\) of semilength
\(n\), all sufficiently large base-\(b\) matching indices are exactly
\(\lfloor b^j/E_q\rfloor\), one for each sufficiently large \(j\).
MacMahon's polynomial \(\genfrac{[}{]}{0pt}{}{2n}{n}_{\!q}/[n+1]_q\) has the corresponding
list \(\lfloor(1-q^{-1})b^j/E_q\rfloor\). These statements use polynomial
evaluation also when \(q\) is not a prime power.
If \(R\) is a finite associative unital ring of order \(b\ge2\), the
number of matching dimensions \(n\le X\) for \(|\mathrm{GL}_n(R)|\)
is \(\log_b X+O_R(\log^* X)\), as it is for singular matrix counts.
To recover the invertible list, write the quotient by the Jacobson
radical as \(R/J(R)=\prod_i M_{d_i}(\mathbb F_{q_i})\), and put
\[
 r_n=\prod_i\prod_{t=1}^{d_i n}(1-q_i^{-t}),\qquad C_R=\prod_i E_{q_i}.
\]
At each sufficiently large scale \(j=n^2-k\) put
\(s=\lceil C_R b^j\rceil,\ \Delta=s-C_Rb^j\), and
\(a_n=b^j(r_n-C_R)\). The list includes \(s-1\) exactly when
\(\Delta>a_{s-1}\) and \(s\) exactly when \(\Delta\le a_s\).
These two candidates exhaust the scale.
For \(R=\mathbb F_q\) in radix \(q\) the two complete
eventual lists are \(\lfloor E_q q^j\rfloor\) and
\(\lfloor(1-E_q)q^j\rfloor\). See
\companion{gaussian-and-catalan-prefixes} and
\companion{matrix-and-parking-prefixes}.

\paragraph{Irreducible polynomials.}
For primes \(p\), let \(U_n\) count monic irreducibles of degree \(n\)
over \(\mathbb F_{p^r}\), \(r\ge1\). In radix \(p^v,\ v\ge1\), the
matching degrees \(n\ge2\) are exactly those with
\[
 n=\lfloor p^m\sqrt p\rfloor,\qquad
 \{p^m\sqrt p\}<\tfrac12,\qquad
 rn-(2m+1)\in v\mathbb Z_{\ge0}\quad(m\ge0).
\]
Degree one matches precisely when \(v\mid r\). For \(p=2,v=1\) the
test reads a zero digit in the binary expansion of \(\sqrt2\).
In this binary case, for every extension degree \(r\ge1\) and every
\(\varepsilon>0\), the number of matches in \([X,X^{1+\varepsilon}]\)
tends to infinity as \(X\to\infty\). In odd
characteristic and radix \(p\), infinitude amounts to infinitely many
digits of \(\sqrt p\) strictly below \((p-1)/2\). Already \(p=3\)
asks about infinitely many ternary zeros in \(\sqrt3\). See
\companion{necklace-prefixes}.

\paragraph{The \(n\)th prime.}
Let \(P_n\) denote the \(n\)th prime and
\(H(y)=\#\{n:yn\le P_n<y(n+1)\}\). For \(0<a<13/30\) and
\(\eta=\exp(-aQ)\), the local mean at every large real \(Q\) satisfies
\[
           \eta^{-1}\int_Q^{Q+\eta}H(y)\,dy\sim Q .
\]
The proof applies the short-interval prime number theorem of Guth and
Maynard~\cite{gm} at the first crossing of \(P_n\ge Qn\). For an
integer \(Q\ge3\), \(H(Q)+1\) is the number of solutions of
\(\pi(Qn-1)=n-1\). The value \(H(16)=0\) illustrates the distinction
between the local mean and its value at a prescribed power. At each sufficiently large power \(Q=b^k\), the short-interval lower bound of Li~\cite{li-primes} yields a growing cluster of indices sharing any fixed fraction below \(0.48\) of their leading digits with their primes. Infinitude
for the unscaled primes in a fixed radix remains unresolved here.
See \companion{nth-prime-prefix-crossings}.

\paragraph{Three fixed sequences and the decimal obstruction.}
For \(B_n\) in radix four the branch multiplier is \(16\), and the
method leaves infinitude open. For the ordinary Catalan number
\(B_n/(n+1)\), finiteness in binary is equivalent to all five numbers
\(\pi^{-1/5}4^{r/5}\), \(0\le r<5\), eventually using only \(0,1\)
in base four. Finiteness in ternary for \((3n)!/(n!)^3\) is equivalent
to both \(3^{1/4}/\sqrt{2\pi}\) and \(3^{3/4}/\sqrt{2\pi}\) eventually
using only ternary digits \(1,2\) with no adjacent ones. Neither
digit condition has been decided here. See
\companion{walk-and-trinomial-prefixes} and
\companion{radix-entropy-and-walks}.

Finally, in the decimal \(p=2\) model the common multiplier really
can keep both branches away from success. Inside the complementary
semicircle choose a closed arc of width \(w\) satisfying
\(1/(\sqrt{10}-1)<w<1/2\). On the real line the dilation by \(\sqrt{10}\)
of each lifted interval contains an integer translate of the entire
interval, since the range of eligible translates has length greater
than one. Nesting gives \(c>0\) with \(c(\sqrt{10})^j\bmod1\) confined
to the arc for every \(j\ge0\). Equation~\eqref{eq:inverse} then
excludes large matches for the integer part of the model with \(C=c^2\)
and any given correction \(\gamma\). Behavior of individual decimal
walk counts requires more information about their constants.
See \companion{seven-base-priority}.

\begin{thebibliography}{99}
\bibitem{pol-sym}
B. Adamczewski and Y. Bugeaud,
\href{https://doi.org/10.1155/IMRN/2006/75968}{Real and \(p\)-adic expansions involving symmetric patterns},
Int. Math. Res. Not. \textbf{2006}, Art. ID 75968, 17 pp.

\bibitem{rev-complexity}
B.~Adamczewski and Y.~Bugeaud,
\href{https://doi.org/10.4007/annals.2007.165.547}{On the complexity of algebraic numbers I. Expansions in integer bases},
Ann. of Math. (2) \textbf{165} (2007), 547--565.

\bibitem{met-baier-ghosh}
S. Baier and A. Ghosh,
\href{https://doi.org/10.1093/qmath/hau016}{Diophantine approximation on lines with prime constraints},
Q. J. Math. \textbf{66} (2015), 1--12.

\bibitem{met-mtp}
V. Beresnevich and S. Velani,
\href{https://doi.org/10.4007/annals.2006.164.971}{A mass transference principle and the Duffin--Schaeffer conjecture for Hausdorff measures},
Ann. of Math. (2) \textbf{164} (2006), 971--992.

\bibitem{met-borosh-fraenkel}
I. Borosh and A. S. Fraenkel,
\href{https://doi.org/10.1016/1385-7258(72)90055-8}{A generalization of Jarn\'ik's theorem on Diophantine approximations},
Indag. Math. \textbf{34} (1972), 193--201.

\bibitem{app-brent}
R.~P. Brent, \emph{Asymptotic approximation of central binomial
coefficients with rigorous error bounds}, Open J. Math. Sci.
\textbf{5} (2021), 380--386.
\href{https://doi.org/10.30538/oms2021.0173}{doi:10.30538/oms2021.0173}.

\bibitem{be-digits}
Y. Bugeaud and J.-H. Evertse,
\href{https://doi.org/10.4064/aa133-3-3}{On two notions of complexity of algebraic numbers},
Acta Arith. \textbf{133} (2008), 221--250.

\bibitem{met-bugeaud-kim}
Y. Bugeaud and D. H. Kim,
\href{https://doi.org/10.5802/aif.3134}{On the expansions of real numbers in two integer bases},
Ann. Inst. Fourier (Grenoble) \textbf{67} (2017), 2225--2235.

\bibitem{met-bugeaud-liao}
Y. Bugeaud and L. Liao,
\href{https://doi.org/10.1017/etds.2014.66}{Uniform Diophantine approximation related to \(b\)-ary and \(\beta\)-expansions},
Ergodic Theory Dynam. Systems \textbf{36} (2016), 1--22.

\bibitem{zec-chang-miller}
S.~Chang and S.~J.~Miller,
\href{https://arxiv.org/abs/2309.00090}{Benford's law under Zeckendorf expansion},
Fibonacci Quart. \textbf{63} (2025), no.~2, 304--335.

\bibitem{cnt-evertse-ferretti}
J.-H. Evertse and R. G. Ferretti,
\href{https://annals.math.princeton.edu/2013/177-2/p04}{A further improvement of the Quantitative Subspace Theorem},
Ann. of Math. (2) \textbf{177} (2013), no.~2, 513--590.

\bibitem{fang}
Z. Fang,
\href{https://arxiv.org/abs/2607.23662}{Self-referential leading digits of exponential sequences: arithmetic structure and certified search},
arXiv:2607.23662 (2026).

\bibitem{app-glendinning-sidorov}
P.~Glendinning and N.~Sidorov, \emph{The doubling map with asymmetrical
holes}, Ergodic Theory Dynam. Systems \textbf{35} (2015), no.~4,
1208--1228.
\href{https://doi.org/10.1017/etds.2013.98}{doi:10.1017/etds.2013.98}.
Numbered results are cited from \href{https://arxiv.org/abs/1302.2486v2}
{arXiv:1302.2486v2}.

\bibitem{gm}
L. Guth and J. Maynard,
\href{https://doi.org/10.4007/annals.2026.203.2.6}{New large value estimates for Dirichlet polynomials}, Ann. of Math. (2) \textbf{203} (2026), 623--675.

\bibitem{met-huang-liu}
J.-J. Huang and J. J. Liu,
\href{https://doi.org/10.1093/imrn/rnz190}{Simultaneous approximation on affine subspaces},
Int. Math. Res. Not. IMRN (2021), 14905--14921.

\bibitem{met-huang}
J.-J. Huang,
\href{https://doi.org/10.1007/s00039-024-00665-y}{Extremal affine subspaces and Khintchine--Jarn\'ik type theorems},
Geom. Funct. Anal. \textbf{34} (2024), 113--163.

\bibitem{met-flatness}
R. Kannan and L. Lov\'asz,
\href{https://doi.org/10.2307/1971436}{Covering minima and lattice-point-free convex bodies},
Ann. of Math. (2) \textbf{128} (1988), 577--602.

\bibitem{li-primes}
R. Li, \href{https://arxiv.org/abs/2308.04458v8}{The number of primes in short intervals and numerical calculations for Harman's sieve}, arXiv:2308.04458v8 (2025).

\bibitem{met-li-li-wu}
B. Li, R. Li, and Y.-F. Wu,
\href{https://doi.org/10.1088/1361-6544/add833}{Multiplicative Diophantine approximation with restricted denominators},
Nonlinearity \textbf{38} (2025), 065012.

\bibitem{met-li-liao}
B. Li, L. Liao, S. Velani, and E. Zorin,
\href{https://doi.org/10.1016/j.aim.2023.108994}{The shrinking target problem for matrix transformations of tori: revisiting the standard problem},
Adv. Math. \textbf{421} (2023), 108994.

\bibitem{app-lindemann}
F.~Lindemann, \emph{\"Uber die Zahl \(\pi\)}, Math. Ann. \textbf{20}
(1882), no.~2, 213--225.
\href{https://doi.org/10.1007/BF01446522}{doi:10.1007/BF01446522}.

\bibitem{pol-mahler}
K. Mahler,
\href{https://doi.org/10.1017/S0004972700021316}{Some suggestions for further research},
Bull. Austral. Math. Soc. \textbf{29} (1984), 101--108.

\bibitem{app-mcintosh}
R.~J. McIntosh, \emph{An asymptotic formula for binomial sums},
J. Number Theory \textbf{58} (1996), no.~1, 158--172.
\href{https://doi.org/10.1006/jnth.1996.0072}{doi:10.1006/jnth.1996.0072}.

\bibitem{met-neukirch}
J. Neukirch,
\href{https://doi.org/10.1007/978-3-662-03983-0}{Algebraic Number Theory},
Grundlehren der mathematischen Wissenschaften, vol.~322, Springer, Berlin, 1999.

\bibitem{pol-ridout}
D. Ridout,
\href{https://doi.org/10.1112/S0025579300001182}{Rational approximations to algebraic numbers},
Mathematika \textbf{4} (1957), 125--131.

\bibitem{zec-northshield}
S.~Northshield,
\href{https://www.mathstat.dal.ca/FQ/Papers1/60-5/northshield.pdf}{Some generalizations of a formula of Reznick},
Fibonacci Quart. \textbf{60} (2022), no.~5, 299--315.

\bibitem{met-rynne}
B. P. Rynne,
\href{https://doi.org/10.1112/S0024609398004536}{Hausdorff dimension and generalized simultaneous Diophantine approximation},
Bull. London Math. Soc. \textbf{30} (1998), 365--376.

\bibitem{met-schlickewei}
H. P. Schlickewei,
\href{https://gdz.sub.uni-goettingen.de/download/pdf/PPN365956996_0018/PPN365956996_0018.pdf}{Linearformen mit algebraischen Koeffizienten},
Manuscripta Math. \textbf{18} (1976), 147--185.

\bibitem{schlickewei}
H. P. Schlickewei,
\href{https://doi.org/10.1016/0022-314X(77)90073-7}{On linear forms with algebraic coefficients and Diophantine equations},
J. Number Theory \textbf{9} (1977), 381--392.

\bibitem{schmidt}
W. M. Schmidt, \href{https://doi.org/10.1007/BF02392334}{Simultaneous approximation to algebraic numbers by rationals}, Acta Math. \textbf{125} (1970), 189--201.

\bibitem{app-straub-zudilin}
A.~Straub and W.~Zudilin, \emph{Sums of powers of binomials, their
Ap\'ery limits, and Franel's suspicions}, Int. Math. Res. Not. IMRN
(2023), no.~11, 9861--9879.
\href{https://doi.org/10.1093/imrn/rnac125}{doi:10.1093/imrn/rnac125}.

\bibitem{zec-stoll}
T.~Stoll,
\href{https://arxiv.org/abs/1009.5360}{Extremal orders of the Zeckendorf sum of digits of powers},
arXiv:1009.5360 (2010).

\bibitem{met-tapia}
L. Tapia,
\href{https://arxiv.org/abs/2602.22512}{Multiplicative Diophantine approximation on planar lines with restricted denominators},
arXiv:2602.22512 (2026).

\bibitem{app-urbanski}
M.~Urba\'nski, \emph{On Hausdorff dimension of invariant sets for
expanding maps of a circle}, Ergodic Theory Dynam. Systems \textbf{6}
(1986), no.~2, 295--309.
\href{https://doi.org/10.1017/S0143385700003461}{doi:10.1017/S0143385700003461}.

\bibitem{vdl}
J. van de Lune,
\href{https://ir.cwi.nl/pub/7417/}{A note on a problem of Erd\H{o}s},
Stichting Mathematisch Centrum, ZN 87/78, Amsterdam, 1978.

\bibitem{app-wagner}
S.~Wagner, \emph{Asymptotics of generalised trinomial coefficients},
\href{https://arxiv.org/abs/1205.5402v3}{arXiv:1205.5402v3}, 2012.

\bibitem{met-wu}
Y.-F. Wu,
\href{https://doi.org/10.1016/j.jmaa.2022.126781}{Inhomogeneous and simultaneous Diophantine approximation in beta dynamical systems},
J. Math. Anal. Appl. \textbf{519} (2023), 126781.

\bibitem{app-zeilberger-zudilin}
D.~Zeilberger and W.~Zudilin, \emph{The irrationality measure of
\(\pi\) is at most \(7.103205334137\ldots\)}, Mosc. J. Comb. Number
Theory \textbf{9} (2020), no.~4, 407--419.
\href{https://doi.org/10.2140/moscow.2020.9.407}{doi:10.2140/moscow.2020.9.407}.

\bibitem{met-zhang-lu}
Y. Zhang and M. L\"u,
\href{https://doi.org/10.1142/S0218348X21501024}{Multiplicative Diophantine approximation of the expansions under different bases on a line},
Fractals \textbf{29} (2021), 2150102.

\bibitem{app-zudilin}
W.~Zudilin, \emph{Two hypergeometric tales and a new irrationality
measure of \(\zeta(2)\)}, Ann. Math. Qu\'ebec \textbf{38} (2014),
no.~1, 101--117.
\href{https://doi.org/10.1007/s40316-014-0016-0}{doi:10.1007/s40316-014-0016-0}.
\end{thebibliography}
\end{document}
